\documentclass[11pt]{article}
\usepackage[T1]{fontenc}
\usepackage[utf8]{inputenc}
\usepackage{amsmath,amssymb,amsthm,mathtools}
\usepackage{geometry,graphicx,booktabs,longtable,tabularx,array}
\usepackage{enumitem,needspace,xcolor,tikz,tcolorbox}
\usetikzlibrary{arrows.meta,positioning,calc,fit,backgrounds,decorations.pathreplacing,matrix,patterns}
\tcbuselibrary{breakable}
\usepackage{caption,hyperref,bookmark,cleveref,microtype}
\usepackage{komorebi-publication}
\renewcommand{\KomorebiTitle}{Komorebi Space}
\renewcommand{\KomorebiSubtitle}{Coherent realizations, obstructions, and descent}
\renewcommand{\KomorebiSeries}{Foundations of the Realization Calculus}
\renewcommand{\KomorebiKind}{Monograph}
\renewcommand{\KomorebiEdition}{Expanded monograph edition 1.4.0 / 11 October 2026}
\renewcommand{\KomorebiShortTitle}{Komorebi Space}
\renewcommand{\KomorebiProfile}{monograph}
\KomorebiSetup
\setcounter{tocdepth}{1}
\setcounter{secnumdepth}{2}
\setlist{nosep,leftmargin=1.7em,topsep=4pt}
\hypersetup{pdfkeywords={Komorebi Space,realization calculus,descent,obstruction,Lean,holonomy,closed range}}
\newtheorem{theorem}{Theorem}[section]
\newtheorem{proposition}[theorem]{Proposition}
\newtheorem{lemma}[theorem]{Lemma}
\newtheorem{corollary}[theorem]{Corollary}
\theoremstyle{definition}
\newtheorem{definition}[theorem]{Definition}
\newtheorem{example}[theorem]{Example}
\theoremstyle{remark}
\newtheorem{remark}[theorem]{Remark}
\numberwithin{equation}{section}
\newcommand{\R}{\mathbb R}
\newcommand{\Z}{\mathbb Z}
\newcommand{\Q}{\mathbb Q}
\newcommand{\N}{\mathbb N}
\newcommand{\K}{\mathbb K}
\newcommand{\V}{\mathcal V}
\newcommand{\W}{\mathcal W}
\newcommand{\id}{\operatorname{id}}
\newcommand{\im}{\operatorname{im}}
\newcommand{\ran}{\operatorname{ran}}
\newcommand{\Hom}{\operatorname{Hom}}
\newcommand{\Ob}{\operatorname{Ob}}
\newcommand{\Ext}{\operatorname{Ext}}
\newcommand{\Match}{\operatorname{Match}}
\newcommand{\Exact}{\operatorname{Exact}}
\newcommand{\dist}{\operatorname{dist}}
\newcommand{\Aut}{\operatorname{Aut}}
\newcommand{\Stab}{\operatorname{Stab}}
\newcommand{\Set}{\mathbf{Set}}
\newcommand{\Cat}{\mathbf{Cat}}
\newcommand{\norm}[1]{\left\lVert#1\right\rVert}
\newcommand{\abs}[1]{\left\lvert#1\right\rvert}
\newcommand{\source}[1]{\par\phantomsection\label{src:#1}{\footnotesize\sffamily\color{KomorebiMuted}Source concordance: \hyperlink{claim:#1}{#1}.}\par\smallskip}
\newcommand{\code}[1]{\texttt{\detokenize{#1}}}
\newcommand{\chap}[1]{\par\Needspace{10\baselineskip}\section{#1}}
\newcommand{\partnote}[2]{\clearpage\part{#1}\noindent\begin{minipage}{.92\linewidth}\large #2\end{minipage}\par\vspace{12pt}}
\newtcolorbox{insight}{breakable,colback=KomorebiSoft!55,colframe=KomorebiPine!50,boxrule=.45pt,arc=1pt,left=9pt,right=9pt,top=7pt,bottom=7pt}
\newtcolorbox{scopebox}{breakable,colback=KomorebiWarm,colframe=KomorebiOchre!55,boxrule=.4pt,arc=1pt,left=9pt,right=9pt,top=7pt,bottom=7pt}
\newcommand{\plate}[3]{\begin{figure}[!htbp]\centering\resizebox{.96\linewidth}{!}{\input{figures/#1}}\caption{#2}\label{#3}\end{figure}}
\tikzset{box/.style={draw=KomorebiPine!70,rounded corners=2pt,fill=KomorebiSoft!50,align=center,inner sep=8pt,font=\small},arr/.style={-{Latex[length=2.2mm]},draw=KomorebiPine,thick},thinarr/.style={-{Latex[length=1.8mm]},draw=KomorebiMuted},lab/.style={font=\small,align=center},dot/.style={circle,fill=KomorebiPine,inner sep=2.3pt}}
\renewcommand{\arraystretch}{1.18}
\raggedbottom
\widowpenalty=7000\clubpenalty=7000
\setlength{\emergencystretch}{3em}

% Compact only the contents entries, without reducing body type size.
\usepackage{etoolbox}
\makeatletter
\patchcmd{\l@section}{1.0em}{0.6em}{}{}
\makeatother

% Generated by checks/edition_stats.py from the current executed reports.
% Generated from current executed reports. Do not hand edit.
\newcommand{\EditionClaims}{80}
\newcommand{\EditionReferences}{211}
\newcommand{\EditionWrittenOnly}{0}
\newcommand{\EditionOriginalChecks}{14,224}
\newcommand{\EditionRecoveredChecks}{2,767}
\newcommand{\EditionRevisionChecks}{17,213}
\newcommand{\EditionExpansionChecks}{96,345}
\newcommand{\EditionDeepeningChecks}{13,817}
\newcommand{\EditionTotalChecks}{144,366}
\newcommand{\EditionFormalModules}{138}
\newcommand{\EditionAuditedDeclarations}{4,745}
\newcommand{\EditionNewTheorems}{54}
\newcommand{\EditionFormalTheorems}{1,507}
\newcommand{\EditionNewConstructions}{17}
\newcommand{\EditionReplays}{1}
\newcommand{\EditionNegativeControls}{7}


% Keep readable half-page figures with their surrounding exposition.
\renewcommand{\topfraction}{.88}
\renewcommand{\bottomfraction}{.80}
\renewcommand{\textfraction}{.10}
\renewcommand{\floatpagefraction}{.78}
\setcounter{topnumber}{2}
\setcounter{bottomnumber}{2}
\setcounter{totalnumber}{3}

\renewcommand{\KomorebiKind}{Lean and evidence guide}
\renewcommand{\KomorebiSubtitle}{Sources, theorem correspondence, and verification}
\renewcommand{\KomorebiEdition}{Companion edition 1.4.0 / 11 October 2026}
\begin{document}
\KomorebiCover
\section*{What this edition establishes}
This guide connects the monograph \emph{Komorebi Space: Coherent realizations, obstructions, and descent} to one integrated Lean project. The canonical source closure joins the foundation, moving-space, and Hilbert holonomy sources. The version 1.2.0 formal baseline closed KS54 through KS58. The preceding expansion developed KS59 through KS70. This edition adds the exposition items KS71 through KS80 from declarations already present in the same accepted source closure. The entire formal tree remains unchanged, and no new Lean compilation is claimed in version 1.4.0.

All \EditionClaims\ principal exposition items now have source correspondences, comprising \EditionReferences\ resolved declaration references. A reference count is not a count of distinct theorems or discoveries. Several printed statements assemble related declarations. Illustrative calculations and explanatory sentences are not automatically promoted to separately formalized claims.

\begin{scopebox}
\textbf{Retained formal baseline: version 1.2.0}\par
The canonical project has \EditionFormalModules\ production modules and one root. The completed clean source replay records, imported axiom audit, positive examples, and deliberate rejection controls are in \path{audit/formal/}. The audit covers \EditionAuditedDeclarations\ project declarations, including compiler-generated declarations. Its permitted axiom set is \code{propext}, \code{Classical.choice}, and \code{Quot.sound}. The pinned external dependency cache is reused rather than rebuilt.
\end{scopebox}

\subsection*{The distinctions that remain important}
A successful compiler run certifies the typed declarations, not a translation of every word of the manuscript. A correct theorem can have narrower hypotheses than an informal description. Its application hypotheses still need to hold in a concrete mathematical model. Formal verification also does not determine historical novelty.

Accordingly, this guide separates retained compiler evidence, imported axiom dependencies, theorem-to-paper correspondence, finite implementation tests, and the remaining mathematical boundaries. The actual Lean types remain authoritative for formal scope.

\subsection*{Files to start with}
The entrypoint is \path{formal/KomorebiSpace.lean}. The canonical sources are in \path{formal/KomorebiSpace/}. The source map is \path{data/claim_ledger.json}, with a CSV companion. The retained formal release summary is \path{audit/formal/release.json}. Each replay has its own complete command receipt, source hashes, axiom report, and logs. The root progress ledger identifies the current restart point.

\clearpage\section{One canonical project, preserved provenance}
\subsection{F: the foundation}
The unchanged archive \path{Komorebi_Space_Lean_v0.20.0.zip} contains 122 production modules and a root. Its historical inventory reports 1,299 named theorem declarations and 696 constructions, including auxiliary results. This tree develops views, partial extension, additive obstruction, tower repair, residual cascades, nonabelian cocycles, effective descent, the composition-site stack, presentation transport, and finite-dimensional amplified holonomy repair.

Its historical strategy included one full cumulative build and a subsequent archive replay reusing a checked 117-module baseline while rebuilding five additions and the root. Those historical records are not relabeled as two clean cumulative replays.

\subsection{M: the focused moving-space closure}
The unchanged archive \path{Komorebi_Moving_Exact_Spaces_v0.32.0.zip} contains 11 production modules and a root. Its historical inventory reports 151 named theorem declarations and 47 constructions. Six modules are inherited dependencies of that focused closure and five are additions. Its recorded verification includes two clean builds of the supplied closure, axiom requests, and negative controls.

The higher version number does not make this archive a replacement for the foundation. It supplies closed-range repair criteria, moving-kernel estimates, constructed perturbed splittings, actual kernel transport, and the quantitative-descent bridge.

\subsection{S: Hilbert holonomy}
The separately recovered \path{HilbertHolonomyRepair.lean} implements amplified holonomy repair with Hilbert fibers and its closed-range necessity criterion. It lay outside both historical archive replay scopes. Version 1.2.0 included it in the canonical source closure and compiled it with its actual dependencies. That accepted evidence is retained here and bound to unchanged source bytes, not borrowed from an unrelated historical report.

\subsection{I: canonical sources and the five version 1.2.0 additions}
The five modules added in version 1.2.0 are \path{FinitePrefixBoundary.lean}, \path{DyadicSupport.lean}, \path{ProtectedHilbertRepair.lean}, \path{ProtectedLinearCertificates.lean}, and \path{FiniteRealizationExistence.lean}. They contain \EditionNewTheorems\ explicit named theorems and lemmas and \EditionNewConstructions\ indexed constructions. These counts include supporting results and do not assert that the classical principles used are newly discovered.

\begin{center}
\begin{tabularx}{\linewidth}{@{}p{.58\linewidth}r@{}}
\toprule Source-selection accounting & Modules\\\midrule
Foundation production sources & 122\\
Focused production sources & 11\\
Separate Hilbert holonomy source & 1\\
Identical shared module counted once & $-1$\\
Version 1.2.0 formalization modules & 5\\\midrule
Canonical production modules & \EditionFormalModules\\
Canonical root & 1\\\bottomrule
\end{tabularx}
\end{center}
The shared module is \code{KernelProjectionBounds}. Its two copies are byte-identical. No inherited production source needed a mathematical or elaboration repair to compile in the canonical union. Both original roots and all frozen archives are preserved separately.

\section{The five formalization gaps that were closed}
\subsection{KS54: finite prefixes do not supply a global point}
The formal model uses the existing \code{Resolution} with fibers \(X_n=\{m\in\N:n\le m\}\) and inclusion restrictions. For every \(N\), the actual partial-section type on the prefix \(0,\ldots,N\) is inhabited by the constant value \(N\). A global point would have one integer value at every coordinate, contradicting the membership condition at its successor.

The proof additionally transfers emptiness to the existing ambient-section type and rejects the existence of a compatible matching family of all prefix sections. This makes the distinction precise: separate prefix witnesses are not matching prefix data.

\subsection{KS55: the dyadic support criterion}
The induction is carried out on the actual \code{DyadicResidual.System}, retaining its changing operator fibers and the starting scale. It proves
\[
 2^N\nmid(k+1)\quad\Longrightarrow\quad (Q_Nr)_k=0.
\]
The corresponding reconstructed coordinate is exact. A nonzero remainder coordinate must lie on the allowed congruence support. The additional valuation form gives exact recovery when \(N>\nu_2(k+1)\).

The induction uses the actual parity-lift identity. It does not replace the varying-fiber system by a scalar model. Neither support theorem asserts that every allowed coordinate is occupied or that coordinatewise recovery implies uniform-norm convergence.

\subsection{KS56: nearest repair with a protected observable}
The source is a complete real Hilbert space. Both the defect \(D\) and protected observable \(p\) are continuous linear maps into the declared normed spaces. Corrections belong to the actual closed kernel \(M=\ker p\), and the operator being studied is \(B=D|_M\).

Feasibility is proved equivalent to the existence of an exact answer preserving \(p(x)\). For feasible \(x\), a raw permitted correction is projected off \(\ker B\) inside \(M\). The resulting displacement is the unique orthogonal permitted correction with image \(Dx\). The actual repaired answer satisfies exactness, preservation, nearest-point uniqueness, and the Pythagorean identity.

No closed-range assumption is needed for nearest-point existence. With a Banach residual target, a uniform residual bound on all feasible inputs is equivalent to closedness of \(\operatorname{ran}B\). The inverse norm is proved to be the least admissible nonnegative constant, including the zero-operator case. Closedness of \(D(E)\) alone is not substituted for this restricted image condition.

\subsection{KS57: literal field-linear certificates}
For linear maps over an arbitrary field, the generic theorem supplies either a protected exact answer or normalized algebraic functionals satisfying the obstruction equations. Algebraic dual separation proves the negative alternative without assuming an elimination routine is correct.

The finite-coordinate specialization proves the literal row-vector statement
\[
 \lambda A=\mu P,\qquad \lambda Ax=1.
\]
It also proves exclusivity and exhaustiveness, identifies the answer fiber as an affine translate of \(\ker A\cap\ker P\), and proves its dimension by rank-nullity for the literal vertically stacked matrix. Zero-dimensional index sets are permitted.

The printed proof uses row reduction, whereas the formal proof uses dual separation. Both establish the stated mathematical alternative. Neither is a universal refinement theorem for the separate Python solver, and neither replaces an integral feasibility problem by a rational one without changing the question.

\subsection{KS58: finite fibers and persistent images}
An explicit functor from \(\N^{op}\) is constructed from the original resolution maps. A coordinate-preserving equivalence identifies its sections with the original \code{Resolution.Point}. The pinned mathlib finite cofiltered-system theorem proves nonemptiness when every fiber is finite and nonempty.

The eventual-range machinery then identifies the range of evaluation at level \(n\) with the intersection of all later restriction images. This proves the persistent-image extension criterion, not surjectivity of evaluation onto every fiber. The printed sequential stabilization proof and the formal cofiltered-system proof have different routes but the same declared conclusion.

\section{Replays, axiom reports, and controls}
\subsection{Pinned runtime of the retained baseline}
For the accepted version 1.2.0 replay, the portable runtime was restored from seven retained offline archive parts. That restoration is recorded evidence, not a new operation in the present editorial expansion. Their payload hashes and CRCs were checked, as was the assembled archive stream. The restored Lean version is 4.34.0-rc2 at commit
\begin{center}\small\path{6a10ac8c22beadecabdbb0919c2b50214762f91d}.\end{center}
The mathlib revision is
\begin{center}\small\path{aeb72f1e9981b6d35b426b16dca7759be3e7ba06}.\end{center}
The verifier checks all nine dependency revisions against the delivered pin manifest. No network operation or floating dependency update is required by the offline replay.

\subsection{Accepted baseline source compilation}
Each completed clean replay begins with an empty local build directory. It resolves the complete local import graph, rejects missing modules and cycles, checks that the root reaches every production file, and compiles each source with one worker, warnings treated as errors, and automatic implicit variables disabled. Every command has an exit code, timing, log hash, and source-bound output record.

The release summary identifies the completed clean replay records. The sources are checked for mutation across each run. A timed-out or failed command is not a pass. Failed development attempts are preserved separately rather than used as final compiler evidence.

\subsection{Imported axiom audit}
The audit enumerates declarations by their defining project module, including generated declarations. It asks Lean for transitive axiom dependencies and rejects any project declaration using an axiom outside \code{propext}, \code{Classical.choice}, and \code{Quot.sound}. Unsafe declarations are rejected as well. The retained environment inventory contains \EditionAuditedDeclarations\ audited project declarations. The lexical source inventory separately counts \EditionFormalTheorems\ explicit theorem and lemma declarations.

The dependency query uses Lean's imported axiom machinery. The compiler, kernel, and imported cached mathematical dependencies remain in the trust boundary. This is not an independent kernel implementation or a full source rebuild of mathlib. Source hashes authenticate the bytes recorded by this package, not an external execution authority.

\subsection{Positive and negative controls}
The positive Lean regression file instantiates protected rational answers, a normalized obstruction, distinct exact answers without the protected constraint, infeasibility under the identity observable, and finite Boolean resolution examples. It also combines separate prefix inhabitance with the proved absence of a global point.

The final control suite has \EditionNegativeControls\ rejection cases. False arithmetic, a proof placeholder, the false finite-prefix implication, universal protected feasibility, an unnormalized certificate, and a wrong preserved value must be rejected by the compiler. A deliberately introduced custom axiom is legal Lean syntax but must be rejected by the dependency policy. That last result is an audit rejection, not a claim that Lean forbids declaring axioms.

\subsection{Exact finite implementation checks}
The independent publication checker executes \EditionOriginalChecks\ assertions. The recovered moving-space checker executes \EditionRecoveredChecks. The revision suite executes \EditionRevisionChecks, including 969 rational certificate problems, independent minor-based rank checks, varying matrix-fiber dyadic examples, and nine implementation rejection controls. The expansion suite adds \EditionExpansionChecks\ exact assertions covering prescribed branches, finite observation, repair mismatch, loop reconstruction, noncommutative boundary comparisons, and uncertainty bounds. The deepening suite adds \EditionDeepeningChecks\ exact assertions for variable-rate paths, closure-boundary approximations, directed costs, and the finite phase extension. Together the five rerun suites contain \EditionTotalChecks\ exact finite assertions. The expansion suite also verifies that all formal files retain their baseline hashes.

These computations check their stated finite models and implementations. The general results identified by principal KS entries come from the formal declarations under their hypotheses. The additional infinite-dimensional example is proved in the manuscript. Neither type of universal conclusion is obtained by extrapolating finite samples. The Python Gaussian proposer remains tested software rather than a universally verified Lean program.

\section{The preceding exposition expansion}
The twelve preceding exposition entries KS59--KS70 did not enlarge the formal source inventory. They expose bridges that were already implemented but insufficiently developed in the reader edition.

\begin{center}
\begin{tabularx}{\linewidth}{@{}p{.18\linewidth}X@{}}
\toprule Entries & Added mathematical development\\\midrule
KS59--KS60 & Stronger prescriptions, quotient naturality, shrinking ambiguity, and a three-view integer choice problem.\\
KS61--KS64 & Scheduled transformations, eventual equality, faithfulness from surjective point projections, the finite-observation image theorem, and cofinal transport.\\
KS65--KS67 & Split exact/residual coordinates, the off-diagonal repair mismatch, its composition law, and bounded-gap residual sampling.\\
KS68 & A concrete failure of literal full-framing invariance, contrasted with marked boundary identifications.\\
KS69--KS70 & Measured residuals, operator uncertainty, and the strict-radius repair theorem when the old kernel is preserved.\\
\bottomrule
\end{tabularx}
\end{center}

The two loop chapters are fully calculated worked applications. The first distinguishes changing inconsistent edge data from changing a vertex candidate. The second distinguishes root-based reconstruction from nearest repair in an explicitly declared Euclidean norm. Their finite calculations are tested independently, but the edition does not assign a new Lean theorem to every numerical specialization or exercise.

The source map has \EditionClaims\ principal items and \EditionReferences\ resolved references. All new references occur in the retained axiom inventory and have unchanged source hashes. Printed proofs remain an editorial mathematical layer. Their presence in the monograph is not a claim of automatic prose-to-Lean translation.

\section{The new convergence and obstruction development}
KS71--KS75 use the retained \path{Exactification}, \path{LinearExactification}, \path{Profiles}, and \path{RepairExamples} modules. The key logical separation is between residual decay, control of total movement, and a comparison between the limits of different paths. The printed examples instantiate these distinctions with exact scalar coordinates. They are not newly compiled instantiations.

KS76--KS77 use \path{ObstructionFloor}. In a normed additive group, zero distance to the correction subgroup means membership in its closure. The manuscript's diagonal Hilbert-space model proves directly that this closure can be strictly larger than the actual correction image. Its infinite-dimensional nonexistence and escape arguments are written proofs, not conclusions inferred from the finite checker and not a new Lean instance.

KS78--KS79 use \path{WeightedGeometry}. They establish extended directed-distance laws only after an independent arrow cost has been declared. The source's zero-distance setoid does not identify categorical isomorphism classes. The free-path example supplies a concrete printed counterexample to that identification.

KS80 uses \path{ProjectiveExample}, with an additive phase in \(\mathbf F_2\). The actual retained theorem concerns the cocycle law and the impossibility of a scalar coboundary. Its matrix realization and eight-element extension are fully worked and exhaustively tested finite interpretations. No new representation-theory module or arbitrary higher-cohomology theorem is claimed.

The model-to-evidence ledger is \path{data/deepening_models.json}. It distinguishes the general retained declarations from printed instantiations, finite arithmetic tests, and bounds on an infinite tail. None of these additions alters the accepted formal tree.

\section{Reading theorem-to-paper correspondence}
The stable KS identifiers organize exposition rather than declaration counts. The source resolver matches names, defining modules, and hashes against the retained accepted compiled-environment inventory. Old F, M, and S references retain their original source paths and are also matched to byte-identical canonical copies. I references point directly to canonical sources, including both inherited modules and the five version 1.2.0 additions.

The consequential hypotheses are retained: transition versus composition covers, feasibility versus uniqueness, identified boundaries, closed-range stability, and the residual bounds required for perturbation. These distinctions remain explicit in the monograph and in the actual Lean types.

\clearpage\section{The correspondence catalogue}
The version 1.2.0 formalization review checked the five written additions against their actual types, including the restricted kernel in KS56, normalization and coefficient field in KS57, and finite nonempty fibers in KS58. The inherited correspondence is preserved. The expanded correspondence is checked against actual declarations and unchanged source hashes in the same accepted baseline. It does not claim a new compiler run. No automatic semantic equivalence checker for natural-language proofs is claimed.

Each entry records the editorial role, actual declaration, source line, and short source-hash prefix. The JSON ledger supplies complete paths and hashes. Repeated references are intentional where a source supports several parts of the exposition.

\textbf{F} is foundation provenance. \textbf{M} is focused moving-space provenance. \textbf{S} is the recovered Hilbert source. \textbf{I} is a canonical source module. All referenced source files belong to the accepted version 1.2.0 compiled closure.
\pdfbookmark[2]{KS01 / Limit universal property}{guide:KS01}
\par\noindent\begin{minipage}{\linewidth}\subsection*{KS01 / Limit universal property}
Grouped source statements. The universe and smallness conventions are retained.
\par\smallskip\noindent{\small\textbf{F} \path{ViewSections.lean}, line 39.\par\noindent\path{KomorebiSpace.ViewSystem.evaluations_jointly_injective}\par{\sffamily\footnotesize Source SHA-256 prefix: 53bb5c493e5cada2.}}\par
\par\smallskip\noindent{\small\textbf{F} \path{ViewSections.lean}, line 99.\par\noindent\path{KomorebiSpace.ViewSystem.Cone.exists_unique_lift}\par{\sffamily\footnotesize Source SHA-256 prefix: 53bb5c493e5cada2.}}\par
\smallskip\noindent{\color{KomorebiMuted}\rule{\linewidth}{.3pt}}\par\end{minipage}\par\medskip
\pdfbookmark[2]{KS02 / Functorial transport}{guide:KS02}
\par\noindent\begin{minipage}{\linewidth}\subsection*{KS02 / Functorial transport}
Direct coordinate-level transport laws.
\par\smallskip\noindent{\small\textbf{F} \path{ViewSections.lean}, line 73.\par\noindent\path{KomorebiSpace.ViewSystem.Morphism.onAmbient_id}\par{\sffamily\footnotesize Source SHA-256 prefix: 53bb5c493e5cada2.}}\par
\par\smallskip\noindent{\small\textbf{F} \path{ViewSections.lean}, line 75.\par\noindent\path{KomorebiSpace.ViewSystem.Morphism.onAmbient_comp}\par{\sffamily\footnotesize Source SHA-256 prefix: 53bb5c493e5cada2.}}\par
\smallskip\noindent{\color{KomorebiMuted}\rule{\linewidth}{.3pt}}\par\end{minipage}\par\medskip
\pdfbookmark[2]{KS03 / Terminal evaluation}{guide:KS03}
\par\noindent\begin{minipage}{\linewidth}\subsection*{KS03 / Terminal evaluation}
An already-global terminal fiber is recovered, not proved inhabited.
\par\smallskip\noindent{\small\textbf{F} \path{ViewSections.lean}, line 142.\par\noindent\path{KomorebiSpace.ViewSystem.terminalEvaluationEquiv}\par{\sffamily\footnotesize Source SHA-256 prefix: 53bb5c493e5cada2.}}\par
\smallskip\noindent{\color{KomorebiMuted}\rule{\linewidth}{.3pt}}\par\end{minipage}\par\medskip
\pdfbookmark[2]{KS04 / Strict section gluing}{guide:KS04}
\par\noindent\begin{minipage}{\linewidth}\subsection*{KS04 / Strict section gluing}
Requires a transition cover, not just object coverage.
\par\smallskip\noindent{\small\textbf{F} \path{SectionGluing.lean}, line 98.\par\noindent\path{KomorebiSpace.ViewSystem.gluingEquiv}\par{\sffamily\footnotesize Source SHA-256 prefix: cf412d554b19ca22.}}\par
\par\smallskip\noindent{\small\textbf{F} \path{SectionGluing.lean}, line 85.\par\noindent\path{KomorebiSpace.ViewSystem.Matching.exists_unique_glue}\par{\sffamily\footnotesize Source SHA-256 prefix: cf412d554b19ca22.}}\par
\smallskip\noindent{\color{KomorebiMuted}\rule{\linewidth}{.3pt}}\par\end{minipage}\par\medskip
\pdfbookmark[2]{KS05 / Object-only cover failure}{guide:KS05}
\par\noindent\begin{minipage}{\linewidth}\subsection*{KS05 / Object-only cover failure}
The Boolean counterexample is implemented.
\par\smallskip\noindent{\small\textbf{F} \path{ViewBoundaryExamples.lean}, line 36.\par\noindent\path{KomorebiSpace.ViewBoundaryExamples.singleton_matching_not_gluable}\par{\sffamily\footnotesize Source SHA-256 prefix: c094eb03ea6b155f.}}\par
\smallskip\noindent{\color{KomorebiMuted}\rule{\linewidth}{.3pt}}\par\end{minipage}\par\medskip
\pdfbookmark[2]{KS06 / Correspondence fibers}{guide:KS06}
\par\noindent\begin{minipage}{\linewidth}\subsection*{KS06 / Correspondence fibers}
Ordinary set fibers, not homotopy fibers.
\par\smallskip\noindent{\small\textbf{F} \path{ViewCorrespondence.lean}, line 43.\par\noindent\path{KomorebiSpace.ViewSystem.bridgeEquiv}\par{\sffamily\footnotesize Source SHA-256 prefix: a89331eabd7662f8.}}\par
\par\smallskip\noindent{\small\textbf{F} \path{ViewCorrespondence.lean}, line 94.\par\noindent\path{KomorebiSpace.ViewSystem.bridge_exists_iff_matching}\par{\sffamily\footnotesize Source SHA-256 prefix: a89331eabd7662f8.}}\par
\smallskip\noindent{\color{KomorebiMuted}\rule{\linewidth}{.3pt}}\par\end{minipage}\par\medskip
\pdfbookmark[2]{KS07 / Nonempty fibers without a section}{guide:KS07}
\par\noindent\begin{minipage}{\linewidth}\subsection*{KS07 / Nonempty fibers without a section}
The one-object Boolean swap has no global fixed choice.
\par\smallskip\noindent{\small\textbf{F} \path{ViewBoundaryExamples.lean}, line 90.\par\noindent\path{KomorebiSpace.ViewBoundaryExamples.swap_no_ambient}\par{\sffamily\footnotesize Source SHA-256 prefix: c094eb03ea6b155f.}}\par
\par\smallskip\noindent{\small\textbf{F} \path{ViewBoundaryExamples.lean}, line 87.\par\noindent\path{KomorebiSpace.ViewBoundaryExamples.swap_transports_bijective}\par{\sffamily\footnotesize Source SHA-256 prefix: c094eb03ea6b155f.}}\par
\smallskip\noindent{\color{KomorebiMuted}\rule{\linewidth}{.3pt}}\par\end{minipage}\par\medskip
\pdfbookmark[2]{KS08 / Additive view differential}{guide:KS08}
\par\noindent\begin{minipage}{\linewidth}\subsection*{KS08 / Additive view differential}
Uses the actual arrow-indexed differential.
\par\smallskip\noindent{\small\textbf{F} \path{SectionObstruction.lean}, line 58.\par\noindent\path{KomorebiSpace.AdditiveViewSystem.differential_comp}\par{\sffamily\footnotesize Source SHA-256 prefix: 9089dd26b6bf34a0.}}\par
\par\smallskip\noindent{\small\textbf{F} \path{SectionObstruction.lean}, line 64.\par\noindent\path{KomorebiSpace.AdditiveViewSystem.differential_zero_iff}\par{\sffamily\footnotesize Source SHA-256 prefix: 9089dd26b6bf34a0.}}\par
\smallskip\noindent{\color{KomorebiMuted}\rule{\linewidth}{.3pt}}\par\end{minipage}\par\medskip
\pdfbookmark[2]{KS09 / Additive extension criterion}{guide:KS09}
\par\noindent\begin{minipage}{\linewidth}\subsection*{KS09 / Additive extension criterion}
The quotient is defined independently of solvability.
\par\smallskip\noindent{\small\textbf{F} \path{SectionObstruction.lean}, line 126.\par\noindent\path{KomorebiSpace.AdditiveViewSystem.extensionClass_zero_iff}\par{\sffamily\footnotesize Source SHA-256 prefix: 9089dd26b6bf34a0.}}\par
\par\smallskip\noindent{\small\textbf{F} \path{SectionObstruction.lean}, line 142.\par\noindent\path{KomorebiSpace.AdditiveViewSystem.extensionClass_independent}\par{\sffamily\footnotesize Source SHA-256 prefix: 9089dd26b6bf34a0.}}\par
\smallskip\noindent{\color{KomorebiMuted}\rule{\linewidth}{.3pt}}\par\end{minipage}\par\medskip
\pdfbookmark[2]{KS10 / Extension ambiguity torsor}{guide:KS10}
\par\noindent\begin{minipage}{\linewidth}\subsection*{KS10 / Extension ambiguity torsor}
Translation requires one chosen extension.
\par\smallskip\noindent{\small\textbf{F} \path{SectionExtensionFibers.lean}, line 35.\par\noindent\path{KomorebiSpace.AdditiveViewSystem.extensionParameters}\par{\sffamily\footnotesize Source SHA-256 prefix: e422b1f42cdb28a1.}}\par
\par\smallskip\noindent{\small\textbf{F} \path{SectionExtensionFibers.lean}, line 62.\par\noindent\path{KomorebiSpace.AdditiveViewSystem.unique_extension_iff}\par{\sffamily\footnotesize Source SHA-256 prefix: e422b1f42cdb28a1.}}\par
\smallskip\noindent{\color{KomorebiMuted}\rule{\linewidth}{.3pt}}\par\end{minipage}\par\medskip
\pdfbookmark[2]{KS11 / Parity seed example}{guide:KS11}
\par\noindent\begin{minipage}{\linewidth}\subsection*{KS11 / Parity seed example}
The disconnected free-coordinate variant in the text is an additional written consequence, not a separate Lean declaration.
\par\smallskip\noindent{\small\textbf{F} \path{SectionExtensionExample.lean}, line 65.\par\noindent\path{KomorebiSpace.SectionExtensionExample.extensionClass_zero_iff_even}\par{\sffamily\footnotesize Source SHA-256 prefix: 12f3be43a2401a41.}}\par
\par\smallskip\noindent{\small\textbf{F} \path{SectionExtensionExample.lean}, line 79.\par\noindent\path{KomorebiSpace.SectionExtensionExample.odd_seed_obstructed}\par{\sffamily\footnotesize Source SHA-256 prefix: 12f3be43a2401a41.}}\par
\par\smallskip\noindent{\small\textbf{F} \path{SectionExtensionExample.lean}, line 85.\par\noindent\path{KomorebiSpace.SectionExtensionExample.nonzero_seed_extends}\par{\sffamily\footnotesize Source SHA-256 prefix: 12f3be43a2401a41.}}\par
\smallskip\noindent{\color{KomorebiMuted}\rule{\linewidth}{.3pt}}\par\end{minipage}\par\medskip
\pdfbookmark[2]{KS12 / Exactness at partial sections}{guide:KS12}
\par\noindent\begin{minipage}{\linewidth}\subsection*{KS12 / Exactness at partial sections}
No surjectivity onto the whole obstruction quotient is asserted.
\par\smallskip\noindent{\small\textbf{F} \path{SectionExactSequence.lean}, line 108.\par\noindent\path{KomorebiSpace.AdditiveViewSystem.connecting_kernel_eq_range}\par{\sffamily\footnotesize Source SHA-256 prefix: 242e72f1d4e431ee.}}\par
\par\smallskip\noindent{\small\textbf{F} \path{SectionExactSequence.lean}, line 132.\par\noindent\path{KomorebiSpace.AdditiveViewSystem.restriction_surjective_iff_connecting_zero}\par{\sffamily\footnotesize Source SHA-256 prefix: 242e72f1d4e431ee.}}\par
\smallskip\noindent{\color{KomorebiMuted}\rule{\linewidth}{.3pt}}\par\end{minipage}\par\medskip
\pdfbookmark[2]{KS13 / Extension by stages}{guide:KS13}
\par\noindent\begin{minipage}{\linewidth}\subsection*{KS13 / Extension by stages}
The dependent sum retains the intermediate choice.
\par\smallskip\noindent{\small\textbf{F} \path{SectionExtensionStages.lean}, line 28.\par\noindent\path{KomorebiSpace.ViewSystem.extensionStages}\par{\sffamily\footnotesize Source SHA-256 prefix: 8196e1272a7300f7.}}\par
\par\smallskip\noindent{\small\textbf{F} \path{SectionExtensionStages.lean}, line 50.\par\noindent\path{KomorebiSpace.ViewSystem.extension_nonempty_iff_intermediate}\par{\sffamily\footnotesize Source SHA-256 prefix: 8196e1272a7300f7.}}\par
\smallskip\noindent{\color{KomorebiMuted}\rule{\linewidth}{.3pt}}\par\end{minipage}\par\medskip
\pdfbookmark[2]{KS14 / Protected correction}{guide:KS14}
\par\noindent\begin{minipage}{\linewidth}\subsection*{KS14 / Protected correction}
The protected observable is a declared additive map.
\par\smallskip\noindent{\small\textbf{F} \path{Correction.lean}, line 42.\par\noindent\path{KomorebiSpace.CorrectionProblem.flatten_iff}\par{\sffamily\footnotesize Source SHA-256 prefix: 089d9155e0a43f66.}}\par
\par\smallskip\noindent{\small\textbf{F} \path{Correction.lean}, line 56.\par\noindent\path{KomorebiSpace.CorrectionProblem.obstruction_invariant}\par{\sffamily\footnotesize Source SHA-256 prefix: 089d9155e0a43f66.}}\par
\smallskip\noindent{\color{KomorebiMuted}\rule{\linewidth}{.3pt}}\par\end{minipage}\par\medskip
\pdfbookmark[2]{KS15 / Obstruction witness}{guide:KS15}
\par\noindent\begin{minipage}{\linewidth}\subsection*{KS15 / Obstruction witness}
A witness annihilates all admissible correction defects.
\par\smallskip\noindent{\small\textbf{F} \path{Correction.lean}, line 64.\par\noindent\path{KomorebiSpace.CorrectionProblem.nonzero_certificate}\par{\sffamily\footnotesize Source SHA-256 prefix: 089d9155e0a43f66.}}\par
\smallskip\noindent{\color{KomorebiMuted}\rule{\linewidth}{.3pt}}\par\end{minipage}\par\medskip
\pdfbookmark[2]{KS16 / Exact tower maps and realized points}{guide:KS16}
\par\noindent\begin{minipage}{\linewidth}\subsection*{KS16 / Exact tower maps and realized points}
The maps use a fixed sequential schedule. The record instances also supply identity and composition.
\par\smallskip\noindent{\small\textbf{F} \path{Towers.lean}, line 54.\par\noindent\path{KomorebiSpace.Tower.PreHom.exact_comp}\par{\sffamily\footnotesize Source SHA-256 prefix: fdd2819fff6150d2.}}\par
\par\smallskip\noindent{\small\textbf{F} \path{Realization.lean}, line 66.\par\noindent\path{KomorebiSpace.Resolution.Map.onLimit_comp}\par{\sffamily\footnotesize Source SHA-256 prefix: 90e0a22f2f7eebbe.}}\par
\smallskip\noindent{\color{KomorebiMuted}\rule{\linewidth}{.3pt}}\par\end{minipage}\par\medskip
\pdfbookmark[2]{KS17 / Tower differential composition}{guide:KS17}
\par\noindent\begin{minipage}{\linewidth}\subsection*{KS17 / Tower differential composition}
No additional mixed term is present in this differential.
\par\smallskip\noindent{\small\textbf{F} \path{Towers.lean}, line 75.\par\noindent\path{KomorebiSpace.Tower.PreHom.defect_comp}\par{\sffamily\footnotesize Source SHA-256 prefix: fdd2819fff6150d2.}}\par
\par\smallskip\noindent{\small\textbf{F} \path{TowerCorrection.lean}, line 39.\par\noindent\path{KomorebiSpace.TowerCorrection.exact_iff_differential_zero}\par{\sffamily\footnotesize Source SHA-256 prefix: 6710b485edbbb676.}}\par
\smallskip\noindent{\color{KomorebiMuted}\rule{\linewidth}{.3pt}}\par\end{minipage}\par\medskip
\pdfbookmark[2]{KS18 / Cofinal point equivalence}{guide:KS18}
\par\noindent\begin{minipage}{\linewidth}\subsection*{KS18 / Cofinal point equivalence}
The subsequent alternating example is an explanatory specialization, not a statement that residual norms are invariant.
\par\smallskip\noindent{\small\textbf{F} \path{Cofinal.lean}, line 94.\par\noindent\path{KomorebiSpace.Resolution.cofinalEquiv}\par{\sffamily\footnotesize Source SHA-256 prefix: 0306f5a0a7f05e14.}}\par
\smallskip\noindent{\color{KomorebiMuted}\rule{\linewidth}{.3pt}}\par\end{minipage}\par\medskip
\pdfbookmark[2]{KS19 / Geometric operator repair}{guide:KS19}
\par\noindent\begin{minipage}{\linewidth}\subsection*{KS19 / Geometric operator repair}
The correction map and residual contraction are hypotheses.
\par\smallskip\noindent{\small\textbf{F} \path{RepairDynamics.lean}, line 70.\par\noindent\path{KomorebiSpace.RepairDynamics.exactification_from_correction_operator}\par{\sffamily\footnotesize Source SHA-256 prefix: 419b196d29b2f14d.}}\par
\par\smallskip\noindent{\small\textbf{F} \path{RepairProjection.lean}, line 128.\par\noindent\path{KomorebiSpace.RepairDynamics.Scheme.projection_bound}\par{\sffamily\footnotesize Source SHA-256 prefix: 4292822995be802a.}}\par
\par\smallskip\noindent{\small\textbf{F} \path{RepairProjection.lean}, line 132.\par\noindent\path{KomorebiSpace.RepairDynamics.Scheme.projection_idempotent}\par{\sffamily\footnotesize Source SHA-256 prefix: 4292822995be802a.}}\par
\smallskip\noindent{\color{KomorebiMuted}\rule{\linewidth}{.3pt}}\par\end{minipage}\par\medskip
\pdfbookmark[2]{KS20 / Constructed split tower correction}{guide:KS20}
\par\noindent\begin{minipage}{\linewidth}\subsection*{KS20 / Constructed split tower correction}
Target splittings and the strict propagation gap are retained.
\par\smallskip\noindent{\small\textbf{F} \path{SplitTowerRepair.lean}, line 154.\par\noindent\path{KomorebiSpace.BanachTower.Splitting.differential_correction}\par{\sffamily\footnotesize Source SHA-256 prefix: 2dca3052b0338605.}}\par
\par\smallskip\noindent{\small\textbf{F} \path{SplitTowerRepair.lean}, line 199.\par\noindent\path{KomorebiSpace.BanachTower.Splitting.repair_error}\par{\sffamily\footnotesize Source SHA-256 prefix: 2dca3052b0338605.}}\par
\smallskip\noindent{\color{KomorebiMuted}\rule{\linewidth}{.3pt}}\par\end{minipage}\par\medskip
\pdfbookmark[2]{KS21 / Split semantic and residual coordinates}{guide:KS21}
\par\noindent\begin{minipage}{\linewidth}\subsection*{KS21 / Split semantic and residual coordinates}
A continuous linear equivalence, not an isometry or an absolute semantic coordinate system.
\par\smallskip\noindent{\small\textbf{F} \path{SplitTowerCoordinates.lean}, line 98.\par\noindent\path{KomorebiSpace.BanachTower.Splitting.coordinates}\par{\sffamily\footnotesize Source SHA-256 prefix: c028e9936850deb9.}}\par
\par\smallskip\noindent{\small\textbf{F} \path{SplitTowerCoordinates.lean}, line 68.\par\noindent\path{KomorebiSpace.BanachTower.Splitting.coordinate_inverse}\par{\sffamily\footnotesize Source SHA-256 prefix: c028e9936850deb9.}}\par
\par\smallskip\noindent{\small\textbf{F} \path{SplitTowerSemantics.lean}, line 141.\par\noindent\path{KomorebiSpace.BanachTower.Splitting.residual_observable_injective}\par{\sffamily\footnotesize Source SHA-256 prefix: ecb64d5ac3d5369d.}}\par
\smallskip\noindent{\color{KomorebiMuted}\rule{\linewidth}{.3pt}}\par\end{minipage}\par\medskip
\pdfbookmark[2]{KS22 / Growing identity repair}{guide:KS22}
\par\noindent\begin{minipage}{\linewidth}\subsection*{KS22 / Growing identity repair}
The source bonding maps have the factor one half.
\par\smallskip\noindent{\small\textbf{F} \path{GrowingTower.lean}, line 195.\par\noindent\path{KomorebiSpace.GrowingTower.identityRepair_formula}\par{\sffamily\footnotesize Source SHA-256 prefix: 988f6ae0da7405cd.}}\par
\smallskip\noindent{\color{KomorebiMuted}\rule{\linewidth}{.3pt}}\par\end{minipage}\par\medskip
\pdfbookmark[2]{KS23 / A bounded residual without bounded primitive}{guide:KS23}
\par\noindent\begin{minipage}{\linewidth}\subsection*{KS23 / A bounded residual without bounded primitive}
This is a counterexample at the threshold, not necessity of the gap in every tower.
\par\smallskip\noindent{\small\textbf{F} \path{IdentityTowerObstruction.lean}, line 55.\par\noindent\path{KomorebiSpace.IdentityTowerObstruction.no_bounded_primitive}\par{\sffamily\footnotesize Source SHA-256 prefix: 026f9504c3c35507.}}\par
\par\smallskip\noindent{\small\textbf{F} \path{IdentityTowerObstruction.lean}, line 76.\par\noindent\path{KomorebiSpace.IdentityTowerObstruction.repair_exists_but_right_inverse_does_not}\par{\sffamily\footnotesize Source SHA-256 prefix: 026f9504c3c35507.}}\par
\smallskip\noindent{\color{KomorebiMuted}\rule{\linewidth}{.3pt}}\par\end{minipage}\par\medskip
\pdfbookmark[2]{KS24 / Dependence on target splitting}{guide:KS24}
\par\noindent\begin{minipage}{\linewidth}\subsection*{KS24 / Dependence on target splitting}
The two first-level matrices in the exposition are derived from the displayed source recurrence.
\par\smallskip\noindent{\small\textbf{F} \path{SplittingDependence.lean}, line 101.\par\noindent\path{KomorebiSpace.SplittingDependence.repair_operators_differ}\par{\sffamily\footnotesize Source SHA-256 prefix: b5aa369adbd67f9f.}}\par
\par\smallskip\noindent{\small\textbf{F} \path{SplittingDependence.lean}, line 112.\par\noindent\path{KomorebiSpace.SplittingDependence.protected_profiles_differ}\par{\sffamily\footnotesize Source SHA-256 prefix: b5aa369adbd67f9f.}}\par
\smallskip\noindent{\color{KomorebiMuted}\rule{\linewidth}{.3pt}}\par\end{minipage}\par\medskip
\pdfbookmark[2]{KS25 / Lossless dyadic residual split}{guide:KS25}
\par\noindent\begin{minipage}{\linewidth}\subsection*{KS25 / Lossless dyadic residual split}
Maps are the actual typed dyadic coarsener and its supplied-section lift.
\par\smallskip\noindent{\small\textbf{F} \path{DyadicResidualSplit.lean}, line 39.\par\noindent\path{KomorebiSpace.DyadicResidual.coarsen_pair}\par{\sffamily\footnotesize Source SHA-256 prefix: c4d40d1c7797e5ad.}}\par
\par\smallskip\noindent{\small\textbf{F} \path{DyadicResidualSplit.lean}, line 113.\par\noindent\path{KomorebiSpace.DyadicResidual.coarsen_lift}\par{\sffamily\footnotesize Source SHA-256 prefix: c4d40d1c7797e5ad.}}\par
\par\smallskip\noindent{\small\textbf{F} \path{DyadicResidualSplit.lean}, line 153.\par\noindent\path{KomorebiSpace.DyadicResidual.reconstruction}\par{\sffamily\footnotesize Source SHA-256 prefix: c4d40d1c7797e5ad.}}\par
\smallskip\noindent{\color{KomorebiMuted}\rule{\linewidth}{.3pt}}\par\end{minipage}\par\medskip
\pdfbookmark[2]{KS26 / Cascade telescope and persistent part}{guide:KS26}
\par\noindent\begin{minipage}{\linewidth}\subsection*{KS26 / Cascade telescope and persistent part}
Generic split cascade. No norm convergence is inferred.
\par\smallskip\noindent{\small\textbf{F} \path{CascadeProjections.lean}, line 79.\par\noindent\path{KomorebiSpace.CoarseningCascade.remainder_nested_left}\par{\sffamily\footnotesize Source SHA-256 prefix: b2865b41336df871.}}\par
\par\smallskip\noindent{\small\textbf{F} \path{CascadeProjections.lean}, line 84.\par\noindent\path{KomorebiSpace.CoarseningCascade.remainder_nested_right}\par{\sffamily\footnotesize Source SHA-256 prefix: b2865b41336df871.}}\par
\par\smallskip\noindent{\small\textbf{F} \path{CascadeProjections.lean}, line 136.\par\noindent\path{KomorebiSpace.CoarseningCascade.persistent_iff_details_zero}\par{\sffamily\footnotesize Source SHA-256 prefix: b2865b41336df871.}}\par
\par\smallskip\noindent{\small\textbf{F} \path{CoarseningCascade.lean}, line 73.\par\noindent\path{KomorebiSpace.CoarseningCascade.reconstruction}\par{\sffamily\footnotesize Source SHA-256 prefix: 2719d5a3accb4dde.}}\par
\smallskip\noindent{\color{KomorebiMuted}\rule{\linewidth}{.3pt}}\par\end{minipage}\par\medskip
\pdfbookmark[2]{KS27 / Recovery by density}{guide:KS27}
\par\noindent\begin{minipage}{\linewidth}\subsection*{KS27 / Recovery by density}
Uniform boundedness of the restored remainders is a hypothesis.
\par\smallskip\noindent{\small\textbf{F} \path{CascadeDensity.lean}, line 75.\par\noindent\path{KomorebiSpace.CoarseningCascade.recoverable_iff_closure_finiteCore}\par{\sffamily\footnotesize Source SHA-256 prefix: 11d90a5995354d99.}}\par
\par\smallskip\noindent{\small\textbf{F} \path{CascadeDensity.lean}, line 81.\par\noindent\path{KomorebiSpace.CoarseningCascade.all_recoverable_iff_dense}\par{\sffamily\footnotesize Source SHA-256 prefix: 11d90a5995354d99.}}\par
\smallskip\noindent{\color{KomorebiMuted}\rule{\linewidth}{.3pt}}\par\end{minipage}\par\medskip
\pdfbookmark[2]{KS28 / Dyadic finite-coordinate recovery}{guide:KS28}
\par\noindent\begin{minipage}{\linewidth}\subsection*{KS28 / Dyadic finite-coordinate recovery}
The printed proof is now complete and uniform in the starting scale. The source retains its original dependent-index induction.
\par\smallskip\noindent{\small\textbf{F} \path{DyadicInfiniteReconstruction.lean}, line 71.\par\noindent\path{KomorebiSpace.DyadicResidual.System.infinite_details_injective}\par{\sffamily\footnotesize Source SHA-256 prefix: 0c81242c3c6c0715.}}\par
\par\smallskip\noindent{\small\textbf{F} \path{DyadicInfiniteReconstruction.lean}, line 201.\par\noindent\path{KomorebiSpace.DyadicResidual.System.reconstruction_coordinate}\par{\sffamily\footnotesize Source SHA-256 prefix: 0c81242c3c6c0715.}}\par
\smallskip\noindent{\color{KomorebiMuted}\rule{\linewidth}{.3pt}}\par\end{minipage}\par\medskip
\pdfbookmark[2]{KS29 / Shift cascade and c0}{guide:KS29}
\par\noindent\begin{minipage}{\linewidth}\subsection*{KS29 / Shift cascade and c0}
The scalar expanding-lift boundary is separately implemented in CascadeBoundaryExamples.
\par\smallskip\noindent{\small\textbf{F} \path{SequenceCascade.lean}, line 140.\par\noindent\path{KomorebiSpace.SequenceCascade.recoverable_iff_vanishes}\par{\sffamily\footnotesize Source SHA-256 prefix: 0a384fa460f07336.}}\par
\par\smallskip\noindent{\small\textbf{F} \path{SequenceCascade.lean}, line 194.\par\noindent\path{KomorebiSpace.SequenceCascade.ones_remainder_norm}\par{\sffamily\footnotesize Source SHA-256 prefix: 0a384fa460f07336.}}\par
\par\smallskip\noindent{\small\textbf{F} \path{SequenceCascade.lean}, line 270.\par\noindent\path{KomorebiSpace.SequenceCascade.geometric_reconstruction_error}\par{\sffamily\footnotesize Source SHA-256 prefix: 0a384fa460f07336.}}\par
\smallskip\noindent{\color{KomorebiMuted}\rule{\linewidth}{.3pt}}\par\end{minipage}\par\medskip
\pdfbookmark[2]{KS30 / Gauge groupoid and section transport}{guide:KS30}
\par\noindent\begin{minipage}{\linewidth}\subsection*{KS30 / Gauge groupoid and section transport}
The action and groupoid instances are also read in this module. Composition order is retained.
\par\smallskip\noindent{\small\textbf{F} \path{GroupViewSystems.lean}, line 107.\par\noindent\path{KomorebiSpace.GroupViewSystem.twisted_coherent_iff}\par{\sffamily\footnotesize Source SHA-256 prefix: d2dd13e5f5cc2d78.}}\par
\par\smallskip\noindent{\small\textbf{F} \path{CocycleGaugeGroupoid.lean}, line 95.\par\noindent\path{KomorebiSpace.GroupViewSystem.gaugeSetoid}\par{\sffamily\footnotesize Source SHA-256 prefix: 93393322cd38df5e.}}\par
\smallskip\noindent{\color{KomorebiMuted}\rule{\linewidth}{.3pt}}\par\end{minipage}\par\medskip
\pdfbookmark[2]{KS31 / Nonabelian extension obstruction}{guide:KS31}
\par\noindent\begin{minipage}{\linewidth}\subsection*{KS31 / Nonabelian extension obstruction}
A pointed based-gauge orbit set, not an assumed cohomology group or a general algorithm.
\par\smallskip\noindent{\small\textbf{F} \path{NonabelianExtension.lean}, line 96.\par\noindent\path{KomorebiSpace.GroupViewSystem.extensionOrbit_neutral_iff}\par{\sffamily\footnotesize Source SHA-256 prefix: 7834ff4193ad9a05.}}\par
\par\smallskip\noindent{\small\textbf{F} \path{NonabelianExtension.lean}, line 114.\par\noindent\path{KomorebiSpace.GroupViewSystem.extensionOrbit_independent}\par{\sffamily\footnotesize Source SHA-256 prefix: 7834ff4193ad9a05.}}\par
\smallskip\noindent{\color{KomorebiMuted}\rule{\linewidth}{.3pt}}\par\end{minipage}\par\medskip
\pdfbookmark[2]{KS32 / Nonabelian extension torsor}{guide:KS32}
\par\noindent\begin{minipage}{\linewidth}\subsection*{KS32 / Nonabelian extension torsor}
Parameters act by right multiplication on a chosen extension.
\par\smallskip\noindent{\small\textbf{F} \path{NonabelianExtensionFibers.lean}, line 51.\par\noindent\path{KomorebiSpace.GroupViewSystem.extensionParameters}\par{\sffamily\footnotesize Source SHA-256 prefix: 486c0bbabc34525c.}}\par
\par\smallskip\noindent{\small\textbf{F} \path{NonabelianExtensionFibers.lean}, line 107.\par\noindent\path{KomorebiSpace.GroupViewSystem.unique_extension_iff}\par{\sffamily\footnotesize Source SHA-256 prefix: 486c0bbabc34525c.}}\par
\smallskip\noindent{\color{KomorebiMuted}\rule{\linewidth}{.3pt}}\par\end{minipage}\par\medskip
\pdfbookmark[2]{KS33 / S3 loop model}{guide:KS33}
\par\noindent\begin{minipage}{\linewidth}\subsection*{KS33 / S3 loop model}
Cocycle coherence is not the existence of a twisted global section.
\par\smallskip\noindent{\small\textbf{F} \path{CocycleExamples.lean}, line 205.\par\noindent\path{KomorebiSpace.CocycleExamples.s3_orbit_dichotomy}\par{\sffamily\footnotesize Source SHA-256 prefix: 6761e0443a1c1ad7.}}\par
\par\smallskip\noindent{\small\textbf{F} \path{CocycleExamples.lean}, line 169.\par\noindent\path{KomorebiSpace.CocycleExamples.actual_automorphisms_noncommute}\par{\sffamily\footnotesize Source SHA-256 prefix: 6761e0443a1c1ad7.}}\par
\par\smallskip\noindent{\small\textbf{F} \path{CocycleExamples.lean}, line 103.\par\noindent\path{KomorebiSpace.CocycleExamples.involution_trivial_iff}\par{\sffamily\footnotesize Source SHA-256 prefix: 6761e0443a1c1ad7.}}\par
\smallskip\noindent{\color{KomorebiMuted}\rule{\linewidth}{.3pt}}\par\end{minipage}\par\medskip
\pdfbookmark[2]{KS34 / Effective groupoid descent}{guide:KS34}
\par\noindent\begin{minipage}{\linewidth}\subsection*{KS34 / Effective groupoid descent}
Unframed gauge groupoids and composition-covered full patches. Uniqueness respects the specified identifications.
\par\smallskip\noindent{\small\textbf{F} \path{DescentEquivalence.lean}, line 104.\par\noindent\path{KomorebiSpace.GroupViewSystem.CocycleDescentData.descentEquivalence}\par{\sffamily\footnotesize Source SHA-256 prefix: c7ca73d617a578ed.}}\par
\par\smallskip\noindent{\small\textbf{F} \path{DescentComparison.lean}, line 79.\par\noindent\path{KomorebiSpace.GroupViewSystem.CocycleDescentData.Effectuation.exists_unique_comparison}\par{\sffamily\footnotesize Source SHA-256 prefix: afc5bf933284a648.}}\par
\smallskip\noindent{\color{KomorebiMuted}\rule{\linewidth}{.3pt}}\par\end{minipage}\par\medskip
\pdfbookmark[2]{KS35 / Composition topology}{guide:KS35}
\par\noindent\begin{minipage}{\linewidth}\subsection*{KS35 / Composition topology}
An actual Grothendieck topology, with the standard generated-sieve characterization.
\par\smallskip\noindent{\small\textbf{F} \path{CompositionSite.lean}, line 57.\par\noindent\path{KomorebiSpace.CompositionSite.topology}\par{\sffamily\footnotesize Source SHA-256 prefix: ac958295633fc50f.}}\par
\par\smallskip\noindent{\small\textbf{F} \path{CompositionSite.lean}, line 143.\par\noindent\path{KomorebiSpace.CompositionSite.ofArrows_covering_iff}\par{\sffamily\footnotesize Source SHA-256 prefix: ac958295633fc50f.}}\par
\smallskip\noindent{\color{KomorebiMuted}\rule{\linewidth}{.3pt}}\par\end{minipage}\par\medskip
\pdfbookmark[2]{KS36 / Standard stackhood}{guide:KS36}
\par\noindent\begin{minipage}{\linewidth}\subsection*{KS36 / Standard stackhood}
Target is mathlib Pseudofunctor.IsStack on the chosen composition site.
\par\smallskip\noindent{\small\textbf{F} \path{CompositionStack.lean}, line 178.\par\noindent\path{KomorebiSpace.GroupViewSystem.patch_isStack}\par{\sffamily\footnotesize Source SHA-256 prefix: 656bbbc1ddaba844.}}\par
\smallskip\noindent{\color{KomorebiMuted}\rule{\linewidth}{.3pt}}\par\end{minipage}\par\medskip
\pdfbookmark[2]{KS37 / Equivalent view presentations}{guide:KS37}
\par\noindent\begin{minipage}{\linewidth}\subsection*{KS37 / Equivalent view presentations}
Prescribed domains are isomorphism closed and coefficient systems are really pulled back.
\par\smallskip\noindent{\small\textbf{F} \path{SectionViewEquivalence.lean}, line 63.\par\noindent\path{KomorebiSpace.ViewSystem.viewEquivalence}\par{\sffamily\footnotesize Source SHA-256 prefix: 999bb5c576cefc7b.}}\par
\par\smallskip\noindent{\small\textbf{F} \path{RelativeSections.lean}, line 66.\par\noindent\path{KomorebiSpace.ViewSystem.relativeViewEquivalence}\par{\sffamily\footnotesize Source SHA-256 prefix: c28de5ab91b69828.}}\par
\par\smallskip\noindent{\small\textbf{F} \path{RelativeAmbiguity.lean}, line 43.\par\noindent\path{KomorebiSpace.GroupViewSystem.invisibleViewEquivalence}\par{\sffamily\footnotesize Source SHA-256 prefix: a24937a0816305f9.}}\par
\smallskip\noindent{\color{KomorebiMuted}\rule{\linewidth}{.3pt}}\par\end{minipage}\par\medskip
\pdfbookmark[2]{KS38 / Identified-boundary realization groupoid}{guide:KS38}
\par\noindent\begin{minipage}{\linewidth}\subsection*{KS38 / Identified-boundary realization groupoid}
Groupoid homotopy fiber construction. No new higher-category interface is asserted.
\par\smallskip\noindent{\small\textbf{F} \path{MarkedRealization.lean}, line 64.\par\noindent\path{KomorebiSpace.MarkedRealization.nonempty_iff}\par{\sffamily\footnotesize Source SHA-256 prefix: a3ad06525f28defd.}}\par
\par\smallskip\noindent{\small\textbf{F} \path{MarkedRealization.lean}, line 73.\par\noindent\path{KomorebiSpace.MarkedRealization.endomorphism_condition}\par{\sffamily\footnotesize Source SHA-256 prefix: a3ad06525f28defd.}}\par
\smallskip\noindent{\color{KomorebiMuted}\rule{\linewidth}{.3pt}}\par\end{minipage}\par\medskip
\pdfbookmark[2]{KS39 / Marked transport across a square}{guide:KS39}
\par\noindent\begin{minipage}{\linewidth}\subsection*{KS39 / Marked transport across a square}
The square commutes through a specified natural isomorphism and the horizontal functors are equivalences.
\par\smallskip\noindent{\small\textbf{F} \path{MarkedTransport.lean}, line 117.\par\noindent\path{KomorebiSpace.MarkedRealization.transportEquivalence}\par{\sffamily\footnotesize Source SHA-256 prefix: 99fc0cf03cfdbdf9.}}\par
\par\smallskip\noindent{\small\textbf{F} \path{MarkedTransport.lean}, line 128.\par\noindent\path{KomorebiSpace.MarkedRealization.transportHomEquiv}\par{\sffamily\footnotesize Source SHA-256 prefix: 99fc0cf03cfdbdf9.}}\par
\smallskip\noindent{\color{KomorebiMuted}\rule{\linewidth}{.3pt}}\par\end{minipage}\par\medskip
\pdfbookmark[2]{KS40 / Transported path telescope}{guide:KS40}
\par\noindent\begin{minipage}{\linewidth}\subsection*{KS40 / Transported path telescope}
Known bounded invertible transports. Path gain is not replaced by length without contractivity.
\par\smallskip\noindent{\small\textbf{F} \path{AmplifiedTransport.lean}, line 151.\par\noindent\path{KomorebiSpace.AmplifiedTransport.Connection.walkSum_defect}\par{\sffamily\footnotesize Source SHA-256 prefix: da79596a2f437a07.}}\par
\par\smallskip\noindent{\small\textbf{F} \path{AmplifiedTransport.lean}, line 239.\par\noindent\path{KomorebiSpace.AmplifiedTransport.Connection.walk_error}\par{\sffamily\footnotesize Source SHA-256 prefix: da79596a2f437a07.}}\par
\smallskip\noindent{\color{KomorebiMuted}\rule{\linewidth}{.3pt}}\par\end{minipage}\par\medskip
\pdfbookmark[2]{KS41 / Holonomy classification}{guide:KS41}
\par\noindent\begin{minipage}{\linewidth}\subsection*{KS41 / Holonomy classification}
A bounded linear equivalence. It is an isometry only in the isometric subclass.
\par\smallskip\noindent{\small\textbf{F} \path{AmplifiedHolonomy.lean}, line 174.\par\noindent\path{KomorebiSpace.AmplifiedTransport.Rooted.parameterEquiv}\par{\sffamily\footnotesize Source SHA-256 prefix: f46a0ba9482bc51c.}}\par
\par\smallskip\noindent{\small\textbf{F} \path{AmplifiedHolonomy.lean}, line 104.\par\noindent\path{KomorebiSpace.AmplifiedTransport.Rooted.fixedSpace_route_independent}\par{\sffamily\footnotesize Source SHA-256 prefix: f46a0ba9482bc51c.}}\par
\smallskip\noindent{\color{KomorebiMuted}\rule{\linewidth}{.3pt}}\par\end{minipage}\par\medskip
\pdfbookmark[2]{KS42 / Orthogonal exact repair}{guide:KS42}
\par\noindent\begin{minipage}{\linewidth}\subsection*{KS42 / Orthogonal exact repair}
Requires the declared Hilbert geometry, not closed residual range.
\par\smallskip\noindent{\small\textbf{M} \path{HilbertKernelProjection.lean}, line 245.\par\noindent\path{KomorebiSpace.HilbertKernelProjection.unconditional_exact_repair}\par{\sffamily\footnotesize Source SHA-256 prefix: 5a30a4de9c18ed34.}}\par
\par\smallskip\noindent{\small\textbf{M} \path{HilbertKernelProjection.lean}, line 120.\par\noindent\path{KomorebiSpace.HilbertKernelProjection.projection_best}\par{\sffamily\footnotesize Source SHA-256 prefix: 5a30a4de9c18ed34.}}\par
\par\smallskip\noindent{\small\textbf{M} \path{HilbertKernelProjection.lean}, line 134.\par\noindent\path{KomorebiSpace.HilbertKernelProjection.projection_unique}\par{\sffamily\footnotesize Source SHA-256 prefix: 5a30a4de9c18ed34.}}\par
\smallskip\noindent{\color{KomorebiMuted}\rule{\linewidth}{.3pt}}\par\end{minipage}\par\medskip
\pdfbookmark[2]{KS43 / Closed range equivalence}{guide:KS43}
\par\noindent\begin{minipage}{\linewidth}\subsection*{KS43 / Closed range equivalence}
Banach target in the stated theorem. The explicit infinite diagonal illustration is written, not newly compiled.
\par\smallskip\noindent{\small\textbf{M} \path{HilbertKernelProjection.lean}, line 210.\par\noindent\path{KomorebiSpace.HilbertKernelProjection.closed_range_iff_bound}\par{\sffamily\footnotesize Source SHA-256 prefix: 5a30a4de9c18ed34.}}\par
\par\smallskip\noindent{\small\textbf{M} \path{HilbertKernelProjection.lean}, line 219.\par\noindent\path{KomorebiSpace.HilbertKernelProjection.closed_range_iff_near_exact}\par{\sffamily\footnotesize Source SHA-256 prefix: 5a30a4de9c18ed34.}}\par
\par\smallskip\noindent{\small\textbf{M} \path{HilbertKernelProjection.lean}, line 163.\par\noindent\path{KomorebiSpace.HilbertKernelProjection.gain_characterization}\par{\sffamily\footnotesize Source SHA-256 prefix: 5a30a4de9c18ed34.}}\par
\smallskip\noindent{\color{KomorebiMuted}\rule{\linewidth}{.3pt}}\par\end{minipage}\par\medskip
\pdfbookmark[2]{KS44 / Amplified holonomy gain}{guide:KS44}
\par\noindent\begin{minipage}{\linewidth}\subsection*{KS44 / Amplified holonomy gain}
The finite-dimensional result retains foundation provenance. The Hilbert necessity theorem retains separate-source provenance. Both are now included in the canonical compiler and axiom record.
\par\smallskip\noindent{\small\textbf{F} \path{AmplifiedRepair.lean}, line 204.\par\noindent\path{KomorebiSpace.AmplifiedTransport.Rooted.orthogonalRepair_error}\par{\sffamily\footnotesize Source SHA-256 prefix: e4009d19b75ab9cf.}}\par
\par\smallskip\noindent{\small\textbf{S} \path{HilbertHolonomyRepair.lean}, line 194.\par\noindent\path{KomorebiSpace.AmplifiedTransport.Rooted.Hilbert.closed_range_iff_near_parallel}\par{\sffamily\footnotesize Source SHA-256 prefix: 5baf9b3a40b4fdd8.}}\par
\smallskip\noindent{\color{KomorebiMuted}\rule{\linewidth}{.3pt}}\par\end{minipage}\par\medskip
\pdfbookmark[2]{KS45 / Flat transport amplification}{guide:KS45}
\par\noindent\begin{minipage}{\linewidth}\subsection*{KS45 / Flat transport amplification}
Both directed edge orientations use the original maximum norm.
\par\smallskip\noindent{\small\textbf{F} \path{AmplifiedExamples.lean}, line 139.\par\noindent\path{KomorebiSpace.AmplifiedTransport.Examples.exponential_amplification}\par{\sffamily\footnotesize Source SHA-256 prefix: e95779e1b481e79e.}}\par
\par\smallskip\noindent{\small\textbf{F} \path{AmplifiedExamples.lean}, line 148.\par\noindent\path{KomorebiSpace.AmplifiedTransport.Examples.length_only_bound_false}\par{\sffamily\footnotesize Source SHA-256 prefix: e95779e1b481e79e.}}\par
\smallskip\noindent{\color{KomorebiMuted}\rule{\linewidth}{.3pt}}\par\end{minipage}\par\medskip
\pdfbookmark[2]{KS46 / Two-stage joint projection}{guide:KS46}
\par\noindent\begin{minipage}{\linewidth}\subsection*{KS46 / Two-stage joint projection}
Second projection is inside the locally exact subspace.
\par\smallskip\noindent{\small\textbf{M} \path{QuantitativeDescent.lean}, line 61.\par\noindent\path{KomorebiSpace.QuantitativeDescent.repair_eq_projection}\par{\sffamily\footnotesize Source SHA-256 prefix: 0a35307bbdb5fdc3.}}\par
\par\smallskip\noindent{\small\textbf{M} \path{QuantitativeDescent.lean}, line 121.\par\noindent\path{KomorebiSpace.QuantitativeDescent.error_squared}\par{\sffamily\footnotesize Source SHA-256 prefix: 0a35307bbdb5fdc3.}}\par
\par\smallskip\noindent{\small\textbf{M} \path{QuantitativeDescent.lean}, line 165.\par\noindent\path{KomorebiSpace.QuantitativeDescent.repair_error_squared}\par{\sffamily\footnotesize Source SHA-256 prefix: 0a35307bbdb5fdc3.}}\par
\par\smallskip\noindent{\small\textbf{M} \path{QuantitativeDescent.lean}, line 133.\par\noindent\path{KomorebiSpace.QuantitativeDescent.operator_interaction}\par{\sffamily\footnotesize Source SHA-256 prefix: 0a35307bbdb5fdc3.}}\par
\smallskip\noindent{\color{KomorebiMuted}\rule{\linewidth}{.3pt}}\par\end{minipage}\par\medskip
\pdfbookmark[2]{KS47 / Joint closed range}{guide:KS47}
\par\noindent\begin{minipage}{\linewidth}\subsection*{KS47 / Joint closed range}
Closed local range is retained. The printed bound is a sufficient, not necessarily optimal, gain.
\par\smallskip\noindent{\small\textbf{M} \path{DescentClosedRange.lean}, line 93.\par\noindent\path{KomorebiSpace.QuantitativeDescent.joint_closed_iff_matching_closed}\par{\sffamily\footnotesize Source SHA-256 prefix: 394f9a40603290cd.}}\par
\par\smallskip\noindent{\small\textbf{M} \path{DescentClosedRange.lean}, line 70.\par\noindent\path{KomorebiSpace.QuantitativeDescent.stable_repair}\par{\sffamily\footnotesize Source SHA-256 prefix: 394f9a40603290cd.}}\par
\smallskip\noindent{\color{KomorebiMuted}\rule{\linewidth}{.3pt}}\par\end{minipage}\par\medskip
\pdfbookmark[2]{KS48 / Two-sided projector motion}{guide:KS48}
\par\noindent\begin{minipage}{\linewidth}\subsection*{KS48 / Two-sided projector motion}
Both independent residual certificates are supplied. No inclusion of kernels is assumed.
\par\smallskip\noindent{\small\textbf{M} \path{MovingKernelBounds.lean}, line 95.\par\noindent\path{KomorebiSpace.MovingKernel.projection_motion}\par{\sffamily\footnotesize Source SHA-256 prefix: 0edf4cbd9a1b1063.}}\par
\par\smallskip\noindent{\small\textbf{M} \path{MovingKernelBounds.lean}, line 164.\par\noindent\path{KomorebiSpace.MovingKernel.projection_lipschitz}\par{\sffamily\footnotesize Source SHA-256 prefix: 0edf4cbd9a1b1063.}}\par
\smallskip\noindent{\color{KomorebiMuted}\rule{\linewidth}{.3pt}}\par\end{minipage}\par\medskip
\pdfbookmark[2]{KS49 / Derived perturbed certificate}{guide:KS49}
\par\noindent\begin{minipage}{\linewidth}\subsection*{KS49 / Derived perturbed certificate}
Banach target, supplied split surjectivity, and strict a times epsilon less than one.
\par\smallskip\noindent{\small\textbf{M} \path{SplitKernelPerturbation.lean}, line 92.\par\noindent\path{KomorebiSpace.SplitKernel.Certificate.perturbed_right_inverse}\par{\sffamily\footnotesize Source SHA-256 prefix: 1e60514ffacf125c.}}\par
\par\smallskip\noindent{\small\textbf{M} \path{SplitKernelPerturbation.lean}, line 98.\par\noindent\path{KomorebiSpace.SplitKernel.Certificate.perturbed_section_bound}\par{\sffamily\footnotesize Source SHA-256 prefix: 1e60514ffacf125c.}}\par
\par\smallskip\noindent{\small\textbf{M} \path{SplitKernelPerturbation.lean}, line 140.\par\noindent\path{KomorebiSpace.SplitKernel.Certificate.projection_motion}\par{\sffamily\footnotesize Source SHA-256 prefix: 1e60514ffacf125c.}}\par
\smallskip\noindent{\color{KomorebiMuted}\rule{\linewidth}{.3pt}}\par\end{minipage}\par\medskip
\pdfbookmark[2]{KS50 / Transport between moving kernels}{guide:KS50}
\par\noindent\begin{minipage}{\linewidth}\subsection*{KS50 / Transport between moving kernels}
Transport depends on the splitting and is not generally orthogonal projection.
\par\smallskip\noindent{\small\textbf{M} \path{SplitKernelPerturbation.lean}, line 216.\par\noindent\path{KomorebiSpace.SplitKernel.Certificate.kernelEquiv}\par{\sffamily\footnotesize Source SHA-256 prefix: 1e60514ffacf125c.}}\par
\par\smallskip\noindent{\small\textbf{M} \path{SplitKernelPerturbation.lean}, line 226.\par\noindent\path{KomorebiSpace.SplitKernel.Certificate.kernel_transport_error}\par{\sffamily\footnotesize Source SHA-256 prefix: 1e60514ffacf125c.}}\par
\smallskip\noindent{\color{KomorebiMuted}\rule{\linewidth}{.3pt}}\par\end{minipage}\par\medskip
\pdfbookmark[2]{KS51 / Exact moving-line example}{guide:KS51}
\par\noindent\begin{minipage}{\linewidth}\subsection*{KS51 / Exact moving-line example}
The all-real projector formula is separate from the strict-radius certificate interface.
\par\smallskip\noindent{\small\textbf{M} \path{MovingKernelExamples.lean}, line 91.\par\noindent\path{KomorebiSpace.MovingKernelExamples.projection_formula}\par{\sffamily\footnotesize Source SHA-256 prefix: 8ce528290f97bfde.}}\par
\par\smallskip\noindent{\small\textbf{M} \path{MovingKernelExamples.lean}, line 142.\par\noindent\path{KomorebiSpace.MovingKernelExamples.motion_opNorm_exact}\par{\sffamily\footnotesize Source SHA-256 prefix: 8ce528290f97bfde.}}\par
\par\smallskip\noindent{\small\textbf{M} \path{MovingKernelExamples.lean}, line 199.\par\noindent\path{KomorebiSpace.MovingKernelExamples.kernel_transport_formula}\par{\sffamily\footnotesize Source SHA-256 prefix: 8ce528290f97bfde.}}\par
\smallskip\noindent{\color{KomorebiMuted}\rule{\linewidth}{.3pt}}\par\end{minipage}\par\medskip
\pdfbookmark[2]{KS52 / Rank loss and radius endpoint}{guide:KS52}
\par\noindent\begin{minipage}{\linewidth}\subsection*{KS52 / Rank loss and radius endpoint}
The rank-changing diagonal pair is a written elementary example. Source correspondence covers the general loss estimate and the radius boundary.
\par\smallskip\noindent{\small\textbf{M} \path{MovingKernelBounds.lean}, line 148.\par\noindent\path{KomorebiSpace.MovingKernel.lost_exact_unit_forces_gain}\par{\sffamily\footnotesize Source SHA-256 prefix: 0edf4cbd9a1b1063.}}\par
\par\smallskip\noindent{\small\textbf{M} \path{MovingKernelExamples.lean}, line 227.\par\noindent\path{KomorebiSpace.MovingKernelExamples.closed_radius_failure}\par{\sffamily\footnotesize Source SHA-256 prefix: 8ce528290f97bfde.}}\par
\smallskip\noindent{\color{KomorebiMuted}\rule{\linewidth}{.3pt}}\par\end{minipage}\par\medskip
\pdfbookmark[2]{KS53 / Motion of quantitative descent}{guide:KS53}
\par\noindent\begin{minipage}{\linewidth}\subsection*{KS53 / Motion of quantitative descent}
Uses the original two-stage repair through its proved identification with the joint projector.
\par\smallskip\noindent{\small\textbf{M} \path{MovingDescent.lean}, line 52.\par\noindent\path{KomorebiSpace.MovingDescent.repair_motion}\par{\sffamily\footnotesize Source SHA-256 prefix: 6820862031ad7094.}}\par
\smallskip\noindent{\color{KomorebiMuted}\rule{\linewidth}{.3pt}}\par\end{minipage}\par\medskip
\pdfbookmark[2]{KS54 / Finite prefixes without a point}{guide:KS54}
\par\noindent\begin{minipage}{\linewidth}\subsection*{KS54 / Finite prefixes without a point}
Uses the original Resolution, its actual finite partial sections, and its Point/Ambient types. Separately nonempty prefixes do not form a matching family.
\par\smallskip\noindent{\small\textbf{I} \path{FinitePrefixBoundary.lean}, line 28.\par\noindent\path{KomorebiSpace.FinitePrefixBoundary.every_prefix_nonempty}\par{\sffamily\footnotesize Source SHA-256 prefix: 92abbc5b41581de0.}}\par
\par\smallskip\noindent{\small\textbf{I} \path{FinitePrefixBoundary.lean}, line 31.\par\noindent\path{KomorebiSpace.FinitePrefixBoundary.no_point}\par{\sffamily\footnotesize Source SHA-256 prefix: 92abbc5b41581de0.}}\par
\par\smallskip\noindent{\small\textbf{I} \path{FinitePrefixBoundary.lean}, line 40.\par\noindent\path{KomorebiSpace.FinitePrefixBoundary.no_ambient}\par{\sffamily\footnotesize Source SHA-256 prefix: 92abbc5b41581de0.}}\par
\par\smallskip\noindent{\small\textbf{I} \path{FinitePrefixBoundary.lean}, line 44.\par\noindent\path{KomorebiSpace.FinitePrefixBoundary.no_matching_prefixes}\par{\sffamily\footnotesize Source SHA-256 prefix: 92abbc5b41581de0.}}\par
\smallskip\noindent{\color{KomorebiMuted}\rule{\linewidth}{.3pt}}\par\end{minipage}\par\medskip
\pdfbookmark[2]{KS55 / Sharper dyadic support criterion}{guide:KS55}
\par\noindent\begin{minipage}{\linewidth}\subsection*{KS55 / Sharper dyadic support criterion}
Scale-uniform induction on the actual varying-fiber dyadic residual system. Nondivisibility implies exact coordinate recovery; allowed support does not imply occupancy.
\par\smallskip\noindent{\small\textbf{I} \path{DyadicSupport.lean}, line 19.\par\noindent\path{KomorebiSpace.DyadicResidual.System.finite_details_zero_coordinate_of_not_dvd}\par{\sffamily\footnotesize Source SHA-256 prefix: b355d709ce765a28.}}\par
\par\smallskip\noindent{\small\textbf{I} \path{DyadicSupport.lean}, line 62.\par\noindent\path{KomorebiSpace.DyadicResidual.System.remainder_coordinate_zero_of_not_dvd}\par{\sffamily\footnotesize Source SHA-256 prefix: b355d709ce765a28.}}\par
\par\smallskip\noindent{\small\textbf{I} \path{DyadicSupport.lean}, line 71.\par\noindent\path{KomorebiSpace.DyadicResidual.System.reconstruction_coordinate_of_not_dvd}\par{\sffamily\footnotesize Source SHA-256 prefix: b355d709ce765a28.}}\par
\par\smallskip\noindent{\small\textbf{I} \path{DyadicSupport.lean}, line 78.\par\noindent\path{KomorebiSpace.DyadicResidual.System.nonzero_remainder_requires_divisibility}\par{\sffamily\footnotesize Source SHA-256 prefix: b355d709ce765a28.}}\par
\par\smallskip\noindent{\small\textbf{I} \path{DyadicSupport.lean}, line 85.\par\noindent\path{KomorebiSpace.DyadicResidual.System.reconstruction_after_valuation}\par{\sffamily\footnotesize Source SHA-256 prefix: b355d709ce765a28.}}\par
\smallskip\noindent{\color{KomorebiMuted}\rule{\linewidth}{.3pt}}\par\end{minipage}\par\medskip
\pdfbookmark[2]{KS56 / Nearest protected Hilbert repair}{guide:KS56}
\par\noindent\begin{minipage}{\linewidth}\subsection*{KS56 / Nearest protected Hilbert repair}
Actual restriction to ker p, with a real Hilbert source and continuous linear maps. Nearest existence needs feasibility but not closed range. The stability equivalence is stated with a Banach residual target.
\par\smallskip\noindent{\small\textbf{I} \path{ProtectedHilbertRepair.lean}, line 33.\par\noindent\path{KomorebiSpace.ProtectedHilbertRepair.feasible_iff_answer}\par{\sffamily\footnotesize Source SHA-256 prefix: 88ec0dbd0638e624.}}\par
\par\smallskip\noindent{\small\textbf{I} \path{ProtectedHilbertRepair.lean}, line 67.\par\noindent\path{KomorebiSpace.ProtectedHilbertRepair.leastCorrection_unique}\par{\sffamily\footnotesize Source SHA-256 prefix: 88ec0dbd0638e624.}}\par
\par\smallskip\noindent{\small\textbf{I} \path{ProtectedHilbertRepair.lean}, line 143.\par\noindent\path{KomorebiSpace.ProtectedHilbertRepair.exists_unique_nearest}\par{\sffamily\footnotesize Source SHA-256 prefix: 88ec0dbd0638e624.}}\par
\par\smallskip\noindent{\small\textbf{I} \path{ProtectedHilbertRepair.lean}, line 109.\par\noindent\path{KomorebiSpace.ProtectedHilbertRepair.repair_pythagoras}\par{\sffamily\footnotesize Source SHA-256 prefix: 88ec0dbd0638e624.}}\par
\par\smallskip\noindent{\small\textbf{I} \path{ProtectedHilbertRepair.lean}, line 178.\par\noindent\path{KomorebiSpace.ProtectedHilbertRepair.closed_range_iff_protected_bound}\par{\sffamily\footnotesize Source SHA-256 prefix: 88ec0dbd0638e624.}}\par
\par\smallskip\noindent{\small\textbf{I} \path{ProtectedHilbertRepair.lean}, line 204.\par\noindent\path{KomorebiSpace.ProtectedHilbertRepair.gain_characterization}\par{\sffamily\footnotesize Source SHA-256 prefix: 88ec0dbd0638e624.}}\par
\par\smallskip\noindent{\small\textbf{I} \path{ProtectedHilbertRepair.lean}, line 213.\par\noindent\path{KomorebiSpace.ProtectedHilbertRepair.gain_zero_iff}\par{\sffamily\footnotesize Source SHA-256 prefix: 88ec0dbd0638e624.}}\par
\smallskip\noindent{\color{KomorebiMuted}\rule{\linewidth}{.3pt}}\par\end{minipage}\par\medskip
\pdfbookmark[2]{KS57 / Finite positive and negative certificates}{guide:KS57}
\par\noindent\begin{minipage}{\linewidth}\subsection*{KS57 / Finite positive and negative certificates}
The arbitrary-field theorem gives literal matrix row witnesses and affine ambiguity. Its formal proof uses algebraic dual separation, not verification of Gaussian elimination or a refinement proof for the Python implementation.
\par\smallskip\noindent{\small\textbf{I} \path{ProtectedLinearCertificates.lean}, line 155.\par\noindent\path{KomorebiSpace.ProtectedLinearCertificates.matrix_certificate_iff}\par{\sffamily\footnotesize Source SHA-256 prefix: dca61c5b0e63b51e.}}\par
\par\smallskip\noindent{\small\textbf{I} \path{ProtectedLinearCertificates.lean}, line 177.\par\noindent\path{KomorebiSpace.ProtectedLinearCertificates.matrix_exclusive_alternative}\par{\sffamily\footnotesize Source SHA-256 prefix: dca61c5b0e63b51e.}}\par
\par\smallskip\noindent{\small\textbf{I} \path{ProtectedLinearCertificates.lean}, line 112.\par\noindent\path{KomorebiSpace.ProtectedLinearCertificates.answerEquivAmbiguity}\par{\sffamily\footnotesize Source SHA-256 prefix: dca61c5b0e63b51e.}}\par
\par\smallskip\noindent{\small\textbf{I} \path{ProtectedLinearCertificates.lean}, line 204.\par\noindent\path{KomorebiSpace.ProtectedLinearCertificates.matrix_ambiguity_dimension}\par{\sffamily\footnotesize Source SHA-256 prefix: dca61c5b0e63b51e.}}\par
\smallskip\noindent{\color{KomorebiMuted}\rule{\linewidth}{.3pt}}\par\end{minipage}\par\medskip
\pdfbookmark[2]{KS58 / Finite-fiber inverse-limit existence}{guide:KS58}
\par\noindent\begin{minipage}{\linewidth}\subsection*{KS58 / Finite-fiber inverse-limit existence}
An explicit coordinate-preserving equivalence links the original Resolution to the pinned mathlib finite cofiltered-system theorem. Persistent images use actual restriction ranges. Finite nonempty fibers do not imply arbitrary prescribed values extend.
\par\smallskip\noindent{\small\textbf{I} \path{FiniteRealizationExistence.lean}, line 28.\par\noindent\path{KomorebiSpace.Resolution.FiniteExistence.pointEquivSections}\par{\sffamily\footnotesize Source SHA-256 prefix: 8dfc9804c24395cb.}}\par
\par\smallskip\noindent{\small\textbf{I} \path{FiniteRealizationExistence.lean}, line 56.\par\noindent\path{KomorebiSpace.Resolution.FiniteExistence.nonempty_point}\par{\sffamily\footnotesize Source SHA-256 prefix: 8dfc9804c24395cb.}}\par
\par\smallskip\noindent{\small\textbf{I} \path{FiniteRealizationExistence.lean}, line 60.\par\noindent\path{KomorebiSpace.Resolution.FiniteExistence.nonempty_ambient}\par{\sffamily\footnotesize Source SHA-256 prefix: 8dfc9804c24395cb.}}\par
\par\smallskip\noindent{\small\textbf{I} \path{FiniteRealizationExistence.lean}, line 70.\par\noindent\path{KomorebiSpace.Resolution.FiniteExistence.evaluation_range}\par{\sffamily\footnotesize Source SHA-256 prefix: 8dfc9804c24395cb.}}\par
\par\smallskip\noindent{\small\textbf{I} \path{FiniteRealizationExistence.lean}, line 87.\par\noindent\path{KomorebiSpace.Resolution.FiniteExistence.extension_iff_persistent}\par{\sffamily\footnotesize Source SHA-256 prefix: 8dfc9804c24395cb.}}\par
\smallskip\noindent{\color{KomorebiMuted}\rule{\linewidth}{.3pt}}\par\end{minipage}\par\medskip
\pdfbookmark[2]{KS59 / Obstructions under stronger prescription}{guide:KS59}
\par\noindent\begin{minipage}{\linewidth}\subsection*{KS59 / Obstructions under stronger prescription}
Additive view systems and nested prescribed domains. The forgetful quotient maps are surjective, not generally injective. Both quotient naturality and persistence refer to the actual extension class.
\par\smallskip\noindent{\small\textbf{I} \path{ObstructionRestriction.lean}, line 21.\par\noindent\path{KomorebiSpace.AdditiveViewSystem.boundaries_antitone}\par{\sffamily\footnotesize Source SHA-256 prefix: 3a99be42c1d3bbbe.}}\par
\par\smallskip\noindent{\small\textbf{I} \path{ObstructionRestriction.lean}, line 54.\par\noindent\path{KomorebiSpace.AdditiveViewSystem.forgetObstruction_surjective}\par{\sffamily\footnotesize Source SHA-256 prefix: 3a99be42c1d3bbbe.}}\par
\par\smallskip\noindent{\small\textbf{I} \path{ObstructionRestriction.lean}, line 46.\par\noindent\path{KomorebiSpace.AdditiveViewSystem.forgetObstruction_comp}\par{\sffamily\footnotesize Source SHA-256 prefix: 3a99be42c1d3bbbe.}}\par
\par\smallskip\noindent{\small\textbf{I} \path{ObstructionRestriction.lean}, line 61.\par\noindent\path{KomorebiSpace.AdditiveViewSystem.extensionClass_restrict}\par{\sffamily\footnotesize Source SHA-256 prefix: 3a99be42c1d3bbbe.}}\par
\par\smallskip\noindent{\small\textbf{I} \path{ObstructionRestriction.lean}, line 71.\par\noindent\path{KomorebiSpace.AdditiveViewSystem.obstruction_persists_under_prescription}\par{\sffamily\footnotesize Source SHA-256 prefix: 3a99be42c1d3bbbe.}}\par
\smallskip\noindent{\color{KomorebiMuted}\rule{\linewidth}{.3pt}}\par\end{minipage}\par\medskip
\pdfbookmark[2]{KS60 / Prescribed choices and obstructed intermediate steps}{guide:KS60}
\par\noindent\begin{minipage}{\linewidth}\subsection*{KS60 / Prescribed choices and obstructed intermediate steps}
The concrete three-view integer model retains its original restriction maps. The initial extension fiber is parametrized by integers, while a prescribed middle value extends precisely when even.
\par\smallskip\noindent{\small\textbf{I} \path{ExtensionChoiceExample.lean}, line 166.\par\noindent\path{KomorebiSpace.ExtensionChoiceExample.initialExtensionsEquiv}\par{\sffamily\footnotesize Source SHA-256 prefix: 02e8b9542fe55cbc.}}\par
\par\smallskip\noindent{\small\textbf{I} \path{ExtensionChoiceExample.lean}, line 110.\par\noindent\path{KomorebiSpace.ExtensionChoiceExample.middle_obstruction_iff_even}\par{\sffamily\footnotesize Source SHA-256 prefix: 02e8b9542fe55cbc.}}\par
\par\smallskip\noindent{\small\textbf{I} \path{ExtensionChoiceExample.lean}, line 184.\par\noindent\path{KomorebiSpace.ExtensionChoiceExample.nonzero_seed_dead_end}\par{\sffamily\footnotesize Source SHA-256 prefix: 02e8b9542fe55cbc.}}\par
\smallskip\noindent{\color{KomorebiMuted}\rule{\linewidth}{.3pt}}\par\end{minipage}\par\medskip
\pdfbookmark[2]{KS61 / Scheduled transports and eventual equality}{guide:KS61}
\par\noindent\begin{minipage}{\linewidth}\subsection*{KS61 / Scheduled transports and eventual equality}
Set-valued sequential resolutions, monotone schedules, and finite-stage eventual equality. The realized action respects the quotient and composition. This is not an identification with every general pro-category.
\par\smallskip\noindent{\small\textbf{I} \path{ReindexedTransports.lean}, line 96.\par\noindent\path{KomorebiSpace.Resolution.ReindexedMap.comp_congr}\par{\sffamily\footnotesize Source SHA-256 prefix: eb3472ee08977cea.}}\par
\par\smallskip\noindent{\small\textbf{I} \path{ReindexedTransports.lean}, line 139.\par\noindent\path{KomorebiSpace.Resolution.ReindexedMap.onLimit_congr}\par{\sffamily\footnotesize Source SHA-256 prefix: eb3472ee08977cea.}}\par
\par\smallskip\noindent{\small\textbf{I} \path{ReindexedTransports.lean}, line 133.\par\noindent\path{KomorebiSpace.Resolution.ReindexedMap.onLimit_comp}\par{\sffamily\footnotesize Source SHA-256 prefix: eb3472ee08977cea.}}\par
\par\smallskip\noindent{\small\textbf{I} \path{EventualCategory.lean}, line 93.\par\noindent\path{KomorebiSpace.Resolution.realizationFunctor}\par{\sffamily\footnotesize Source SHA-256 prefix: 8e5b5dd147c99604.}}\par
\smallskip\noindent{\color{KomorebiMuted}\rule{\linewidth}{.3pt}}\par\end{minipage}\par\medskip
\pdfbookmark[2]{KS62 / Faithfulness from observable finite stages}{guide:KS62}
\par\noindent\begin{minipage}{\linewidth}\subsection*{KS62 / Faithfulness from observable finite stages}
Faithfulness assumes surjectivity of every source projection from the limit. Surjective bonding maps give a sufficient condition. Without it equality on global points can conceal a finite-stage difference.
\par\smallskip\noindent{\small\textbf{I} \path{RealizationFaithfulness.lean}, line 39.\par\noindent\path{KomorebiSpace.Resolution.ReindexedMap.eventuallyEq_iff_onLimit_eq}\par{\sffamily\footnotesize Source SHA-256 prefix: f8527cf96665e665.}}\par
\par\smallskip\noindent{\small\textbf{I} \path{RealizationFaithfulness.lean}, line 110.\par\noindent\path{KomorebiSpace.Resolution.projectionSurjective_of_surjective_bonding}\par{\sffamily\footnotesize Source SHA-256 prefix: f8527cf96665e665.}}\par
\smallskip\noindent{\color{KomorebiMuted}\rule{\linewidth}{.3pt}}\par\end{minipage}\par\medskip
\pdfbookmark[2]{KS63 / Finite observation characterizes realized transport}{guide:KS63}
\par\noindent\begin{minipage}{\linewidth}\subsection*{KS63 / Finite observation characterizes realized transport}
Every source limit projection is surjective. For each target coordinate one finite source stage must work uniformly for all source points. The theorem does not assert that arbitrary functions have finite observation.
\par\smallskip\noindent{\small\textbf{I} \path{FiniteObservation.lean}, line 92.\par\noindent\path{KomorebiSpace.Resolution.realization_image_iff}\par{\sffamily\footnotesize Source SHA-256 prefix: ad7f071a6cc05534.}}\par
\par\smallskip\noindent{\small\textbf{I} \path{FiniteObservation.lean}, line 101.\par\noindent\path{KomorebiSpace.Resolution.finiteObservationEquiv}\par{\sffamily\footnotesize Source SHA-256 prefix: ad7f071a6cc05534.}}\par
\par\smallskip\noindent{\small\textbf{I} \path{FiniteObservation.lean}, line 64.\par\noindent\path{KomorebiSpace.Resolution.observedPresentation_onLimit}\par{\sffamily\footnotesize Source SHA-256 prefix: ad7f071a6cc05534.}}\par
\smallskip\noindent{\color{KomorebiMuted}\rule{\linewidth}{.3pt}}\par\end{minipage}\par\medskip
\pdfbookmark[2]{KS64 / Cofinal transport equivalence}{guide:KS64}
\par\noindent\begin{minipage}{\linewidth}\subsection*{KS64 / Cofinal transport equivalence}
Increasing cofinal schedules give explicit isomorphisms in the eventual transport category and Hom equivalences. This does not assert equality of quantitative residual norms.
\par\smallskip\noindent{\small\textbf{I} \path{CofinalTransports.lean}, line 80.\par\noindent\path{KomorebiSpace.Resolution.cofinalIso}\par{\sffamily\footnotesize Source SHA-256 prefix: 78e7fe892911c3ac.}}\par
\par\smallskip\noindent{\small\textbf{I} \path{CofinalTransports.lean}, line 112.\par\noindent\path{KomorebiSpace.Resolution.cofinalHomEquiv}\par{\sffamily\footnotesize Source SHA-256 prefix: 78e7fe892911c3ac.}}\par
\par\smallskip\noindent{\small\textbf{I} \path{CofinalTransports.lean}, line 123.\par\noindent\path{KomorebiSpace.Resolution.cofinalHomEquiv_onLimit}\par{\sffamily\footnotesize Source SHA-256 prefix: 78e7fe892911c3ac.}}\par
\smallskip\noindent{\color{KomorebiMuted}\rule{\linewidth}{.3pt}}\par\end{minipage}\par\medskip
\pdfbookmark[2]{KS65 / Split repair coordinates and choice dependence}{guide:KS65}
\par\noindent\begin{minipage}{\linewidth}\subsection*{KS65 / Split repair coordinates and choice dependence}
Continuous linear right inverses of a fixed defect operator determine exact/residual coordinates. Projectors with this common kernel image commute exactly when equal, and equality reflects equality of the right inverses.
\par\smallskip\noindent{\small\textbf{I} \path{SplitDefectCharts.lean}, line 169.\par\noindent\path{KomorebiSpace.SplitDefect.Section.chart}\par{\sffamily\footnotesize Source SHA-256 prefix: 6fdf5a5cba4d307a.}}\par
\par\smallskip\noindent{\small\textbf{I} \path{SplitDefectCharts.lean}, line 81.\par\noindent\path{KomorebiSpace.SplitDefect.Section.project_change}\par{\sffamily\footnotesize Source SHA-256 prefix: 6fdf5a5cba4d307a.}}\par
\par\smallskip\noindent{\small\textbf{I} \path{SplitDefectCharts.lean}, line 96.\par\noindent\path{KomorebiSpace.SplitDefect.Section.projects_commute_iff}\par{\sffamily\footnotesize Source SHA-256 prefix: 6fdf5a5cba4d307a.}}\par
\par\smallskip\noindent{\small\textbf{I} \path{SplitDefectCharts.lean}, line 107.\par\noindent\path{KomorebiSpace.SplitDefect.Section.projects_equal_iff_lifts_equal}\par{\sffamily\footnotesize Source SHA-256 prefix: 6fdf5a5cba4d307a.}}\par
\smallskip\noindent{\color{KomorebiMuted}\rule{\linewidth}{.3pt}}\par\end{minipage}\par\medskip
\pdfbookmark[2]{KS66 / Triangular representation change and repair mismatch}{guide:KS66}
\par\noindent\begin{minipage}{\linewidth}\subsection*{KS66 / Triangular representation change and repair mismatch}
The square of defect operators commutes exactly. The mismatch depends on chosen bounded right inverses. Its composition law uses the declared residual map in the second term. Norm convergence requires the supplied bounded maps and vanishing residuals.
\par\smallskip\noindent{\small\textbf{I} \path{ResidualNaturality.lean}, line 123.\par\noindent\path{KomorebiSpace.SplitDefect.DefectMap.chart_formula}\par{\sffamily\footnotesize Source SHA-256 prefix: accf1d105c612e67.}}\par
\par\smallskip\noindent{\small\textbf{I} \path{ResidualNaturality.lean}, line 77.\par\noindent\path{KomorebiSpace.SplitDefect.DefectMap.repair_change}\par{\sffamily\footnotesize Source SHA-256 prefix: accf1d105c612e67.}}\par
\par\smallskip\noindent{\small\textbf{I} \path{ResidualNaturality.lean}, line 105.\par\noindent\path{KomorebiSpace.SplitDefect.DefectMap.repair_natural_iff}\par{\sffamily\footnotesize Source SHA-256 prefix: accf1d105c612e67.}}\par
\par\smallskip\noindent{\small\textbf{I} \path{ResidualNaturality.lean}, line 137.\par\noindent\path{KomorebiSpace.SplitDefect.DefectMap.mismatch_comp}\par{\sffamily\footnotesize Source SHA-256 prefix: accf1d105c612e67.}}\par
\par\smallskip\noindent{\small\textbf{I} \path{ResidualNaturality.lean}, line 145.\par\noindent\path{KomorebiSpace.SplitDefect.DefectMap.repair_discrepancy_tendsto_zero}\par{\sffamily\footnotesize Source SHA-256 prefix: accf1d105c612e67.}}\par
\smallskip\noindent{\color{KomorebiMuted}\rule{\linewidth}{.3pt}}\par\end{minipage}\par\medskip
\pdfbookmark[2]{KS67 / Residual sampling and repair across scales}{guide:KS67}
\par\noindent\begin{minipage}{\linewidth}\subsection*{KS67 / Residual sampling and repair across scales}
Actual operator-valued tower residuals with contractive bonding maps. A uniform bound on sampling gaps is required for the stated residual operator and repair norm estimates, but not for qualitative preservation of exactness.
\par\smallskip\noindent{\small\textbf{I} \path{ResidualSampling.lean}, line 113.\par\noindent\path{KomorebiSpace.BanachTower.resampling_chain}\par{\sffamily\footnotesize Source SHA-256 prefix: 01c4bf790c0ac8e5.}}\par
\par\smallskip\noindent{\small\textbf{I} \path{ResidualSampling.lean}, line 107.\par\noindent\path{KomorebiSpace.BanachTower.residualSampleCLM_opNorm}\par{\sffamily\footnotesize Source SHA-256 prefix: 01c4bf790c0ac8e5.}}\par
\par\smallskip\noindent{\small\textbf{I} \path{ResidualSampling.lean}, line 129.\par\noindent\path{KomorebiSpace.BanachTower.sample_exact}\par{\sffamily\footnotesize Source SHA-256 prefix: 01c4bf790c0ac8e5.}}\par
\par\smallskip\noindent{\small\textbf{I} \path{SamplingRepair.lean}, line 68.\par\noindent\path{KomorebiSpace.BanachTower.sampling_repair_bound}\par{\sffamily\footnotesize Source SHA-256 prefix: 314018aefc244c41.}}\par
\smallskip\noindent{\color{KomorebiMuted}\rule{\linewidth}{.3pt}}\par\end{minipage}\par\medskip
\pdfbookmark[2]{KS68 / Why literal full framing is presentation dependent}{guide:KS68}
\par\noindent\begin{minipage}{\linewidth}\subsection*{KS68 / Why literal full framing is presentation dependent}
The two-copy equivalent-view model compares actual fully framed gauge groupoids. Failure concerns freezing every representative, not the invariant identified-boundary realization construction.
\par\smallskip\noindent{\small\textbf{I} \path{FramingBoundaryExamples.lean}, line 115.\par\noindent\path{KomorebiSpace.FramingBoundaryExamples.full_framing_not_equivalence}\par{\sffamily\footnotesize Source SHA-256 prefix: c55a4b3576aa167a.}}\par
\par\smallskip\noindent{\small\textbf{I} \path{FramingBoundaryExamples.lean}, line 43.\par\noindent\path{KomorebiSpace.GroupViewSystem.fully_framed_equivalent_iff}\par{\sffamily\footnotesize Source SHA-256 prefix: c55a4b3576aa167a.}}\par
\par\smallskip\noindent{\small\textbf{I} \path{FramingBoundaryExamples.lean}, line 75.\par\noindent\path{KomorebiSpace.FramingBoundaryExamples.copyCocycle_pullback}\par{\sffamily\footnotesize Source SHA-256 prefix: c55a4b3576aa167a.}}\par
\smallskip\noindent{\color{KomorebiMuted}\rule{\linewidth}{.3pt}}\par\end{minipage}\par\medskip
\pdfbookmark[2]{KS69 / Measured and uncertain residual certificates}{guide:KS69}
\par\noindent\begin{minipage}{\linewidth}\subsection*{KS69 / Measured and uncertain residual certificates}
A supplied residual bound for the reference operator. Measurement error and operator error have separate terms. Proximity is to the reference exact space, not automatically to the changed kernel.
\par\smallskip\noindent{\small\textbf{I} \path{ResidualPerturbation.lean}, line 33.\par\noindent\path{KomorebiSpace.ResidualRobustness.Bound.measured}\par{\sffamily\footnotesize Source SHA-256 prefix: 5ceeb960fcb6fb63.}}\par
\par\smallskip\noindent{\small\textbf{I} \path{ResidualPerturbation.lean}, line 59.\par\noindent\path{KomorebiSpace.ResidualRobustness.Bound.uncertain_operator_opNorm}\par{\sffamily\footnotesize Source SHA-256 prefix: 5ceeb960fcb6fb63.}}\par
\smallskip\noindent{\color{KomorebiMuted}\rule{\linewidth}{.3pt}}\par\end{minipage}\par\medskip
\pdfbookmark[2]{KS70 / Stable repair when the old exact space is preserved}{guide:KS70}
\par\noindent\begin{minipage}{\linewidth}\subsection*{KS70 / Stable repair when the old exact space is preserved}
The changed operator annihilates the old kernel and the product of reference gain and perturbation radius is strictly less than one. Kernel equality and the new bound are derived, not assumed from norm closeness alone.
\par\smallskip\noindent{\small\textbf{I} \path{ResidualPerturbation.lean}, line 73.\par\noindent\path{KomorebiSpace.ResidualRobustness.Bound.preserved_kernel_bound}\par{\sffamily\footnotesize Source SHA-256 prefix: 5ceeb960fcb6fb63.}}\par
\par\smallskip\noindent{\small\textbf{I} \path{ResidualPerturbation.lean}, line 96.\par\noindent\path{KomorebiSpace.ResidualRobustness.Bound.kernels_equal}\par{\sffamily\footnotesize Source SHA-256 prefix: 5ceeb960fcb6fb63.}}\par
\par\smallskip\noindent{\small\textbf{I} \path{ResidualPerturbation.lean}, line 139.\par\noindent\path{KomorebiSpace.ResidualRobustness.Bound.perturbed_measured}\par{\sffamily\footnotesize Source SHA-256 prefix: 5ceeb960fcb6fb63.}}\par
\smallskip\noindent{\color{KomorebiMuted}\rule{\linewidth}{.3pt}}\par\end{minipage}\par\medskip
\pdfbookmark[2]{KS71 / Exactification from summable movement}{guide:KS71}
\par\noindent\begin{minipage}{\linewidth}\subsection*{KS71 / Exactification from summable movement}
Complete pseudometric source, nonnegative lower semicontinuous residual, summable step bounds and vanishing residual. The monograph specializes to a metric space. No iteration map or uniqueness across different paths is assumed.
\par\smallskip\noindent{\small\textbf{I} \path{Exactification.lean}, line 23.\par\noindent\path{KomorebiSpace.Exactification.residual_zero_of_limit}\par{\sffamily\footnotesize Source SHA-256 prefix: 7039cb91b9563df0.}}\par
\par\smallskip\noindent{\small\textbf{I} \path{Exactification.lean}, line 36.\par\noindent\path{KomorebiSpace.Exactification.summable_exactification}\par{\sffamily\footnotesize Source SHA-256 prefix: 7039cb91b9563df0.}}\par
\smallskip\noindent{\color{KomorebiMuted}\rule{\linewidth}{.3pt}}\par\end{minipage}\par\medskip
\pdfbookmark[2]{KS72 / Variable-rate exactification and posterior bounds}{guide:KS72}
\par\noindent\begin{minipage}{\linewidth}\subsection*{KS72 / Variable-rate exactification and posterior bounds}
The nonnegative factors may exceed one at individual stages. Their prefix product must tend to zero and the weighted movement series must be summable. The geometric estimate is relative to the residual at the restart index.
\par\smallskip\noindent{\small\textbf{I} \path{Exactification.lean}, line 49.\par\noindent\path{KomorebiSpace.Exactification.productRate}\par{\sffamily\footnotesize Source SHA-256 prefix: 7039cb91b9563df0.}}\par
\par\smallskip\noindent{\small\textbf{I} \path{Exactification.lean}, line 61.\par\noindent\path{KomorebiSpace.Exactification.variable_decay}\par{\sffamily\footnotesize Source SHA-256 prefix: 7039cb91b9563df0.}}\par
\par\smallskip\noindent{\small\textbf{I} \path{Exactification.lean}, line 75.\par\noindent\path{KomorebiSpace.Exactification.variable_exactification}\par{\sffamily\footnotesize Source SHA-256 prefix: 7039cb91b9563df0.}}\par
\par\smallskip\noindent{\small\textbf{I} \path{Exactification.lean}, line 111.\par\noindent\path{KomorebiSpace.Exactification.geometric_limit_bound}\par{\sffamily\footnotesize Source SHA-256 prefix: 7039cb91b9563df0.}}\par
\smallskip\noindent{\color{KomorebiMuted}\rule{\linewidth}{.3pt}}\par\end{minipage}\par\medskip
\pdfbookmark[2]{KS73 / Controlled approximation annihilates protected obstruction}{guide:KS73}
\par\noindent\begin{minipage}{\linewidth}\subsection*{KS73 / Controlled approximation annihilates protected obstruction}
A complete normed linear candidate space with continuous linear defect and protected maps. A convergent protected path supplies an actual admissible correction, not merely closure membership of its residual.
\par\smallskip\noindent{\small\textbf{I} \path{LinearExactification.lean}, line 31.\par\noindent\path{KomorebiSpace.LinearExactification.semantic_limit}\par{\sffamily\footnotesize Source SHA-256 prefix: 4a4c8b2227b88b73.}}\par
\par\smallskip\noindent{\small\textbf{I} \path{LinearExactification.lean}, line 42.\par\noindent\path{KomorebiSpace.LinearExactification.summable_semantic_exactification}\par{\sffamily\footnotesize Source SHA-256 prefix: 4a4c8b2227b88b73.}}\par
\par\smallskip\noindent{\small\textbf{I} \path{LinearExactification.lean}, line 67.\par\noindent\path{KomorebiSpace.LinearExactification.controlled_obstruction_vanishes}\par{\sffamily\footnotesize Source SHA-256 prefix: 4a4c8b2227b88b73.}}\par
\smallskip\noindent{\color{KomorebiMuted}\rule{\linewidth}{.3pt}}\par\end{minipage}\par\medskip
\pdfbookmark[2]{KS74 / Path comparison and nonunique protected answers}{guide:KS74}
\par\noindent\begin{minipage}{\linewidth}\subsection*{KS74 / Path comparison and nonunique protected answers}
A vanishing pairwise distance budget identifies the limits of two convergent paths in a metric space. The separate finite model shows that exactness and one nonzero observable alone do not identify answers. The interpolated trajectories are printed examples.
\par\smallskip\noindent{\small\textbf{I} \path{Exactification.lean}, line 168.\par\noindent\path{KomorebiSpace.PathIndependence.limits_equal_of_budget}\par{\sffamily\footnotesize Source SHA-256 prefix: 7039cb91b9563df0.}}\par
\par\smallskip\noindent{\small\textbf{I} \path{RepairExamples.lean}, line 54.\par\noindent\path{KomorebiSpace.RepairExamples.semantics_alone_not_unique}\par{\sffamily\footnotesize Source SHA-256 prefix: 0996b7d0c9007a24.}}\par
\smallskip\noindent{\color{KomorebiMuted}\rule{\linewidth}{.3pt}}\par\end{minipage}\par\medskip
\pdfbookmark[2]{KS75 / What cofinal residual-profile sampling preserves}{guide:KS75}
\par\noindent\begin{minipage}{\linewidth}\subsection*{KS75 / What cofinal residual-profile sampling preserves}
The convergence equivalence assumes a nonnegative antitone profile and monotone unbounded schedule. It does not preserve the numerical rate. The alternating counterexample shows why profile monotonicity matters for the converse.
\par\smallskip\noindent{\small\textbf{I} \path{Profiles.lean}, line 41.\par\noindent\path{KomorebiSpace.Profiles.cofinal_vanish_iff}\par{\sffamily\footnotesize Source SHA-256 prefix: b775fa29880c278d.}}\par
\par\smallskip\noindent{\small\textbf{I} \path{Profiles.lean}, line 59.\par\noindent\path{KomorebiSpace.Profiles.sampled_alternating_vanishes}\par{\sffamily\footnotesize Source SHA-256 prefix: b775fa29880c278d.}}\par
\par\smallskip\noindent{\small\textbf{I} \path{Profiles.lean}, line 64.\par\noindent\path{KomorebiSpace.Profiles.alternating_does_not_vanish}\par{\sffamily\footnotesize Source SHA-256 prefix: b775fa29880c278d.}}\par
\smallskip\noindent{\color{KomorebiMuted}\rule{\linewidth}{.3pt}}\par\end{minipage}\par\medskip
\pdfbookmark[2]{KS76 / Obstruction floors detect closure, not exact membership}{guide:KS76}
\par\noindent\begin{minipage}{\linewidth}\subsection*{KS76 / Obstruction floors detect closure, not exact membership}
Normed additive residual group and an actual additive subgroup of admissible defects. Zero floor is closure membership. Closedness is separately required to recover exact membership. No nearest point is asserted in the general normed group.
\par\smallskip\noindent{\small\textbf{I} \path{ObstructionFloor.lean}, line 25.\par\noindent\path{KomorebiSpace.ObstructionFloor.lower_bound}\par{\sffamily\footnotesize Source SHA-256 prefix: 7f509aff5210bcbc.}}\par
\par\smallskip\noindent{\small\textbf{I} \path{ObstructionFloor.lean}, line 29.\par\noindent\path{KomorebiSpace.ObstructionFloor.zero_iff_closure}\par{\sffamily\footnotesize Source SHA-256 prefix: 7f509aff5210bcbc.}}\par
\par\smallskip\noindent{\small\textbf{I} \path{ObstructionFloor.lean}, line 33.\par\noindent\path{KomorebiSpace.ObstructionFloor.zero_iff_mem_of_closed}\par{\sffamily\footnotesize Source SHA-256 prefix: 7f509aff5210bcbc.}}\par
\smallskip\noindent{\color{KomorebiMuted}\rule{\linewidth}{.3pt}}\par\end{minipage}\par\medskip
\pdfbookmark[2]{KS77 / Persistent floors block all admissible repair sequences}{guide:KS77}
\par\noindent\begin{minipage}{\linewidth}\subsection*{KS77 / Persistent floors block all admissible repair sequences}
The residuals and admissible subgroups may vary by index. An eventually uniform positive lower floor rules out norm convergence to zero for every admissible correction sequence.
\par\smallskip\noindent{\small\textbf{I} \path{ObstructionFloor.lean}, line 39.\par\noindent\path{KomorebiSpace.ObstructionFloor.persistent_obstruction}\par{\sffamily\footnotesize Source SHA-256 prefix: 7f509aff5210bcbc.}}\par
\smallskip\noindent{\color{KomorebiMuted}\rule{\linewidth}{.3pt}}\par\end{minipage}\par\medskip
\pdfbookmark[2]{KS78 / Independent transport costs and directed geometry}{guide:KS78}
\par\noindent\begin{minipage}{\linewidth}\subsection*{KS78 / Independent transport costs and directed geometry}
Extended nonnegative arrow costs with identity zero and composition subadditivity. The source constructs directed and symmetric distance laws and the zero-distance setoid, not a criterion for categorical isomorphism.
\par\smallskip\noindent{\small\textbf{I} \path{WeightedGeometry.lean}, line 29.\par\noindent\path{KomorebiSpace.TransportCost.distance_le_cost}\par{\sffamily\footnotesize Source SHA-256 prefix: 2f927f94ffb32e80.}}\par
\par\smallskip\noindent{\small\textbf{I} \path{WeightedGeometry.lean}, line 32.\par\noindent\path{KomorebiSpace.TransportCost.distance_self}\par{\sffamily\footnotesize Source SHA-256 prefix: 2f927f94ffb32e80.}}\par
\par\smallskip\noindent{\small\textbf{I} \path{WeightedGeometry.lean}, line 36.\par\noindent\path{KomorebiSpace.TransportCost.distance_triangle}\par{\sffamily\footnotesize Source SHA-256 prefix: 2f927f94ffb32e80.}}\par
\par\smallskip\noindent{\small\textbf{I} \path{WeightedGeometry.lean}, line 60.\par\noindent\path{KomorebiSpace.TransportCost.symDistance_triangle}\par{\sffamily\footnotesize Source SHA-256 prefix: 2f927f94ffb32e80.}}\par
\par\smallskip\noindent{\small\textbf{I} \path{WeightedGeometry.lean}, line 72.\par\noindent\path{KomorebiSpace.TransportCost.zeroDistanceSetoid}\par{\sffamily\footnotesize Source SHA-256 prefix: 2f927f94ffb32e80.}}\par
\smallskip\noindent{\color{KomorebiMuted}\rule{\linewidth}{.3pt}}\par\end{minipage}\par\medskip
\pdfbookmark[2]{KS79 / Zero arrows can collapse object distance}{guide:KS79}
\par\noindent\begin{minipage}{\linewidth}\subsection*{KS79 / Zero arrows can collapse object distance}
A category with zero morphisms and a valid transport cost vanishing on each zero arrow. Raw residual norms are not automatically subadditive costs. Semantic restrictions change the admissible category.
\par\smallskip\noindent{\small\textbf{I} \path{WeightedGeometry.lean}, line 42.\par\noindent\path{KomorebiSpace.TransportCost.zero_cost_distance}\par{\sffamily\footnotesize Source SHA-256 prefix: 2f927f94ffb32e80.}}\par
\par\smallskip\noindent{\small\textbf{I} \path{WeightedGeometry.lean}, line 47.\par\noindent\path{KomorebiSpace.TransportCost.collapse_of_zero_arrows}\par{\sffamily\footnotesize Source SHA-256 prefix: 2f927f94ffb32e80.}}\par
\smallskip\noindent{\color{KomorebiMuted}\rule{\linewidth}{.3pt}}\par\end{minipage}\par\medskip
\pdfbookmark[2]{KS80 / An associative phase with no scalar flattening}{guide:KS80}
\par\noindent\begin{minipage}{\linewidth}\subsection*{KS80 / An associative phase with no scalar flattening}
The retained Lean model is the additive F2-valued phase x2*y1 on F2 squared. It proves normalization, coherence and nontriviality under scalar coboundaries. The real matrix lift, complex-scalar observation and eight-element extension are supplied written derivations and finite tests, not new formal declarations.
\par\smallskip\noindent{\small\textbf{I} \path{ProjectiveExample.lean}, line 21.\par\noindent\path{KomorebiSpace.ProjectiveExample.phase_zero_left}\par{\sffamily\footnotesize Source SHA-256 prefix: 408d4ee8346ecb89.}}\par
\par\smallskip\noindent{\small\textbf{I} \path{ProjectiveExample.lean}, line 22.\par\noindent\path{KomorebiSpace.ProjectiveExample.phase_zero_right}\par{\sffamily\footnotesize Source SHA-256 prefix: 408d4ee8346ecb89.}}\par
\par\smallskip\noindent{\small\textbf{I} \path{ProjectiveExample.lean}, line 24.\par\noindent\path{KomorebiSpace.ProjectiveExample.cocycle}\par{\sffamily\footnotesize Source SHA-256 prefix: 408d4ee8346ecb89.}}\par
\par\smallskip\noindent{\small\textbf{I} \path{ProjectiveExample.lean}, line 31.\par\noindent\path{KomorebiSpace.ProjectiveExample.coboundary_symmetric}\par{\sffamily\footnotesize Source SHA-256 prefix: 408d4ee8346ecb89.}}\par
\par\smallskip\noindent{\small\textbf{I} \path{ProjectiveExample.lean}, line 36.\par\noindent\path{KomorebiSpace.ProjectiveExample.phase_asymmetric}\par{\sffamily\footnotesize Source SHA-256 prefix: 408d4ee8346ecb89.}}\par
\par\smallskip\noindent{\small\textbf{I} \path{ProjectiveExample.lean}, line 39.\par\noindent\path{KomorebiSpace.ProjectiveExample.no_scalar_flattening}\par{\sffamily\footnotesize Source SHA-256 prefix: 408d4ee8346ecb89.}}\par
\smallskip\noindent{\color{KomorebiMuted}\rule{\linewidth}{.3pt}}\par\end{minipage}\par\medskip


\clearpage\section{Reproduction and handoff}
\subsection{Verify the delivered payload first}
From the extracted package root, run
\begin{quote}\small\ttfamily
python3 checks/verify\_payload.py
\end{quote}
This checks the manifest and standalone-font exclusion without rewriting the package. It is an integrity check, not a Lean proof check.

\subsection{A new formal replay}
Supply the pinned portable runtime and choose a new label:
\begin{quote}\small\ttfamily
python3 formal/scripts/verify.py \textbackslash\par
\hspace*{1em}--runtime /path/to/portable-runtime \textbackslash\par
\hspace*{1em}--label independent1
\end{quote}
The verifier refuses to overwrite an existing audit or local build directory. Its command log records the result of each compilation, axiom query, regression file, and rejection control. Local compiled objects are disposable and are not required in the source distribution. The runtime and external dependency cache are not embedded in this package.

The \path{lakefile.toml} and \path{lean-toolchain} describe the same project for a conventional Lean setup. The fully recorded replay path is the supplied offline verifier, not an unexecuted claim about every installation method.

\subsection{Regenerate the reader set}
The root \path{build.sh} reruns source integrity, correspondence, finite checks, figure generation, evidence counts, and all three PDF builds. It does not silently execute another Lean replay. It consumes the completed source-bound formal release record. Re-render the resulting PDFs using \path{checks/render_review.py} and inspect the pages before renewing the hash-bound visual-review record.

Rebuilding may update current timestamps and files. Historical archives and replay evidence retain their own dates and scope. A modified source needs a new formal replay and release binding before it can inherit current verification status.

\section{Remaining scope}
Version 1.2.0 closed the identified source-integration and five-item formalization tasks. Version 1.4.0 expands exposition and finite checks without changing that baseline. Universal-view doctrines, higher stacks, unrestricted semantic realization, nonlinear repair, and general arithmetic correspondences remain separate research questions. Neither a new compiler replay nor universal correctness of the Python implementation is claimed here.

The progress ledger records the expansion, failures, checks, and exact restart point. The publication version describes this reader edition, not completion of every proposed Komorebi extension.
\end{document}
